\section{未锚定账本事件的叙事补充：eastern-jin-sixteen-kingdoms}

\label{sec:anchor-batch-two-eastern-jin-sixteen-kingdoms}

以下事件均为此前正文尚未设置导航目标的账本记录。每一锚点置于来源、制度与地方条件的叙事中，避免把年表、政权分类或战报修辞当作社会全貌。

\subsection{制度与地方资源}

\eventanchor{JH-0367-FORMER-QIN-REBELLIONS}{太和二年（367年）·苻双、苻柳等宗室和地方势力反叛，苻坚与王猛平定内乱}账本记录前秦在太和二年（367年）发生“苻双、苻柳等宗室和地方势力反叛，苻坚与王猛平定内乱”。条目把背景概括为苻坚加强中枢控制触动宗室与地方军权，前秦扩张前先遭内部考验，并指出平叛巩固苻坚和王猛权威，使前秦能集中兵力向东扩张。这个节点进入区域社会时，要经由文书、道路、军队、市场、寺院和家庭等不同中介；因此，政治行动的宣布、地方接收和日常生活的改变不能被视作同一时刻完成。

本条依据《晋书·苻坚载记上》；《晋书·王猛传》；《资治通鉴·晋纪》，证据提示为“主要叛乱与前秦获胜可靠；叛乱者目标、联合关系和镇压规模有争议”。这些材料可支持年序、官方决定、战事或特定地点的线索，却难以直接统计所有人的财富、迁徙、信仰和动机。更稳妥的读法是区分诏令与实施、军报与人口、行政称谓与自我认同，并让地方材料的缺环限制解释的范围。

由此可见，该事件不是孤立的年表点，而是资源、信息与权力关系重新配置的一环。不同地区的官员、社群和家庭可能以不同方式承担征发、保护、迁移或宗教活动；现存来源只能照亮其中一部分，不能替沉默者制造统一的经验。

\eventanchor{JH-0369-FANGTOU-DEFEAT}{太和四年（369年）·桓温北伐前燕，在枋头受慕容垂等反击而败，退军途中损失严重}账本记录东晋与前燕在太和四年（369年）发生“桓温北伐前燕，在枋头受慕容垂等反击而败，退军途中损失严重”。条目把背景概括为桓温欲借前燕内部矛盾北进，但长距离补给和多方势力介入限制行动，并指出桓温威望受损，东晋北伐路线暂告挫败，前秦得以观察并利用燕晋竞争。这个节点进入区域社会时，要经由文书、道路、军队、市场、寺院和家庭等不同中介；因此，政治行动的宣布、地方接收和日常生活的改变不能被视作同一时刻完成。

本条依据《晋书·桓温传》；《晋书·慕容垂载记》；《晋书·苻坚载记上》；《资治通鉴·晋纪》，证据提示为“战败与退军可靠；苻坚援燕的作用、军粮问题和伤亡数字有争议”。这些材料可支持年序、官方决定、战事或特定地点的线索，却难以直接统计所有人的财富、迁徙、信仰和动机。更稳妥的读法是区分诏令与实施、军报与人口、行政称谓与自我认同，并让地方材料的缺环限制解释的范围。

由此可见，该事件不是孤立的年表点，而是资源、信息与权力关系重新配置的一环。不同地区的官员、社群和家庭可能以不同方式承担征发、保护、迁移或宗教活动；现存来源只能照亮其中一部分，不能替沉默者制造统一的经验。

\eventanchor{JH-0370-FORMER-QIN-CONQUERS-YAN}{太和五年（370年）·王猛率前秦军攻灭前燕，慕容暐降服，前秦取得河北和中原北部}账本记录前秦与前燕在太和五年（370年）发生“王猛率前秦军攻灭前燕，慕容暐降服，前秦取得河北和中原北部”。条目把背景概括为前燕在枋头后内部权力失衡，前秦已完成关中整军并掌握进攻时机，并指出前秦成为北方霸主，慕容旧部虽被安置却保留了日后复国的社会网络。这个节点进入区域社会时，要经由文书、道路、军队、市场、寺院和家庭等不同中介；因此，政治行动的宣布、地方接收和日常生活的改变不能被视作同一时刻完成。

本条依据《晋书·苻坚载记上》；《晋书·王猛传》；《晋书·慕容暐载记》；《资治通鉴·晋纪》，证据提示为“前燕灭亡、慕容暐降服与王猛统帅可靠；兵力、城池投降次序和地方接受程度有争议”。这些材料可支持年序、官方决定、战事或特定地点的线索，却难以直接统计所有人的财富、迁徙、信仰和动机。更稳妥的读法是区分诏令与实施、军报与人口、行政称谓与自我认同，并让地方材料的缺环限制解释的范围。

由此可见，该事件不是孤立的年表点，而是资源、信息与权力关系重新配置的一环。不同地区的官员、社群和家庭可能以不同方式承担征发、保护、迁移或宗教活动；现存来源只能照亮其中一部分，不能替沉默者制造统一的经验。

\eventanchor{JH-0371-FORMER-QIN-ANNEXES-CHOUCHI}{咸安元年（371年）·前秦攻取仇池，控制陇南山地通道并扩展对氐族地区的影响}账本记录前秦与仇池氐族政权在咸安元年（371年）发生“前秦攻取仇池，控制陇南山地通道并扩展对氐族地区的影响”。条目把背景概括为前秦统一关中后需控制通往汉中和河西的山地通道，仇池内部力量相对薄弱，并指出前秦扩大西南战略纵深，但山地地方社会并未因此完全纳入直接统治。这个节点进入区域社会时，要经由文书、道路、军队、市场、寺院和家庭等不同中介；因此，政治行动的宣布、地方接收和日常生活的改变不能被视作同一时刻完成。

本条依据《晋书·苻坚载记上》；《晋书·杨氏传》；《资治通鉴·晋纪》，证据提示为“征服仇池与政权更替大体可靠；山地控制范围、部众迁徙和地方治理有争议”。这些材料可支持年序、官方决定、战事或特定地点的线索，却难以直接统计所有人的财富、迁徙、信仰和动机。更稳妥的读法是区分诏令与实施、军报与人口、行政称谓与自我认同，并让地方材料的缺环限制解释的范围。

由此可见，该事件不是孤立的年表点，而是资源、信息与权力关系重新配置的一环。不同地区的官员、社群和家庭可能以不同方式承担征发、保护、迁移或宗教活动；现存来源只能照亮其中一部分，不能替沉默者制造统一的经验。

\eventanchor{JH-0372-XIAOWU-EMPEROR}{咸安二年（372年）·简文帝去世，孝武帝司马曜即位，谢安等高门重臣参与辅政}账本记录东晋在咸安二年（372年）发生“简文帝去世，孝武帝司马曜即位，谢安等高门重臣参与辅政”。条目把背景概括为简文帝在桓温压力下即位时间短，继承必须依赖皇族与高门共识，并指出孝武朝的谢安、桓温集团竞争将影响淝水之战前的国家动员。这个节点进入区域社会时，要经由文书、道路、军队、市场、寺院和家庭等不同中介；因此，政治行动的宣布、地方接收和日常生活的改变不能被视作同一时刻完成。

本条依据《晋书·简文帝纪》；《晋书·孝武帝纪》；《晋书·谢安传》，证据提示为“简文帝卒、孝武帝继位和谢安等辅政可靠；遗诏、皇帝年龄与各家族权力分配有争议”。这些材料可支持年序、官方决定、战事或特定地点的线索，却难以直接统计所有人的财富、迁徙、信仰和动机。更稳妥的读法是区分诏令与实施、军报与人口、行政称谓与自我认同，并让地方材料的缺环限制解释的范围。

由此可见，该事件不是孤立的年表点，而是资源、信息与权力关系重新配置的一环。不同地区的官员、社群和家庭可能以不同方式承担征发、保护、迁移或宗教活动；现存来源只能照亮其中一部分，不能替沉默者制造统一的经验。

\eventanchor{JH-0373-FORMER-QIN-TAKES-LIANG-YI}{宁康元年（373年）·前秦军攻取东晋梁、益二州北部州郡，深入汉中与四川北缘}账本记录前秦与东晋在宁康元年（373年）发生“前秦军攻取东晋梁、益二州北部州郡，深入汉中与四川北缘”。条目把背景概括为前秦灭燕后兵力充足，东晋西部防线因桓氏与中央关系紧张而薄弱，并指出前秦获得进攻长江上游和汉中通道的前沿，东晋战略压力增加。这个节点进入区域社会时，要经由文书、道路、军队、市场、寺院和家庭等不同中介；因此，政治行动的宣布、地方接收和日常生活的改变不能被视作同一时刻完成。

本条依据《晋书·孝武帝纪》；《晋书·苻坚载记下》；《资治通鉴·晋纪》，证据提示为“攻取梁益要地可靠；失地范围、地方投降和前秦行政接管程度有争议”。这些材料可支持年序、官方决定、战事或特定地点的线索，却难以直接统计所有人的财富、迁徙、信仰和动机。更稳妥的读法是区分诏令与实施、军报与人口、行政称谓与自我认同，并让地方材料的缺环限制解释的范围。

由此可见，该事件不是孤立的年表点，而是资源、信息与权力关系重新配置的一环。不同地区的官员、社群和家庭可能以不同方式承担征发、保护、迁移或宗教活动；现存来源只能照亮其中一部分，不能替沉默者制造统一的经验。

\eventanchor{JH-0376-FORMER-QIN-DESTROYS-LIANG}{太元元年（376年）·前秦攻灭前凉，张氏河西政权终结，凉州纳入前秦势力}账本记录前秦与前凉在太元元年（376年）发生“前秦攻灭前凉，张氏河西政权终结，凉州纳入前秦势力”。条目把背景概括为前凉长期依托河西走廊维持自治，前秦统一北方后集中兵力西进，并指出前秦取得河西通道和其赋税、军马资源，但西部地方控制仍需持续协商。这个节点进入区域社会时，要经由文书、道路、军队、市场、寺院和家庭等不同中介；因此，政治行动的宣布、地方接收和日常生活的改变不能被视作同一时刻完成。

本条依据《晋书·苻坚载记下》；《晋书·张轨传》；《资治通鉴·晋纪》；《吐鲁番出土文书》，证据提示为“前凉灭亡与张氏降服可靠；河西城镇实际接收、西域关系和人口迁徙有争议”。这些材料可支持年序、官方决定、战事或特定地点的线索，却难以直接统计所有人的财富、迁徙、信仰和动机。更稳妥的读法是区分诏令与实施、军报与人口、行政称谓与自我认同，并让地方材料的缺环限制解释的范围。

由此可见，该事件不是孤立的年表点，而是资源、信息与权力关系重新配置的一环。不同地区的官员、社群和家庭可能以不同方式承担征发、保护、迁移或宗教活动；现存来源只能照亮其中一部分，不能替沉默者制造统一的经验。

\eventanchor{JH-0376-FORMER-QIN-DESTROYS-DAI}{太元元年（376年）·前秦分兵攻灭代国，拓跋什翼犍败亡，拓跋部暂时分散}账本记录前秦与代国在太元元年（376年）发生“前秦分兵攻灭代国，拓跋什翼犍败亡，拓跋部暂时分散”。条目把背景概括为前秦意图消除北部边地独立军事力量，代国内部亦有继承和部众矛盾，并指出拓跋部短暂受挫但保存了复兴网络，386年将重新建立魏国。这个节点进入区域社会时，要经由文书、道路、军队、市场、寺院和家庭等不同中介；因此，政治行动的宣布、地方接收和日常生活的改变不能被视作同一时刻完成。

本条依据《魏书·序纪》；《晋书·苻坚载记下》；《资治通鉴·晋纪》，证据提示为“灭代与什翼犍死亡大体可靠；内部叛乱、前秦军功和拓跋部众去向有争议”。这些材料可支持年序、官方决定、战事或特定地点的线索，却难以直接统计所有人的财富、迁徙、信仰和动机。更稳妥的读法是区分诏令与实施、军报与人口、行政称谓与自我认同，并让地方材料的缺环限制解释的范围。

由此可见，该事件不是孤立的年表点，而是资源、信息与权力关系重新配置的一环。不同地区的官员、社群和家庭可能以不同方式承担征发、保护、迁移或宗教活动；现存来源只能照亮其中一部分，不能替沉默者制造统一的经验。

\eventanchor{JH-0378-XIANGYANG-SIEGE}{太元三年（378年）·苻丕等围攻襄阳，前秦以汉水交通枢纽为目标向东晋施压}账本记录前秦与东晋在太元三年（378年）发生“苻丕等围攻襄阳，前秦以汉水交通枢纽为目标向东晋施压”。条目把背景概括为前秦控制梁益后希望夺取汉水中游，襄阳是连接关中、荆州与江汉的要地，并指出长期围城牵制东晋荆州防务，并为前秦后续南伐建立前进基地。这个节点进入区域社会时，要经由文书、道路、军队、市场、寺院和家庭等不同中介；因此，政治行动的宣布、地方接收和日常生活的改变不能被视作同一时刻完成。

本条依据《晋书·孝武帝纪》；《晋书·苻坚载记下》；《晋书·朱序传》；《资治通鉴·晋纪》，证据提示为“围城开始和前秦主攻方向可靠；围军数量、城内粮食与东晋援军调度有争议”。这些材料可支持年序、官方决定、战事或特定地点的线索，却难以直接统计所有人的财富、迁徙、信仰和动机。更稳妥的读法是区分诏令与实施、军报与人口、行政称谓与自我认同，并让地方材料的缺环限制解释的范围。

由此可见，该事件不是孤立的年表点，而是资源、信息与权力关系重新配置的一环。不同地区的官员、社群和家庭可能以不同方式承担征发、保护、迁移或宗教活动；现存来源只能照亮其中一部分，不能替沉默者制造统一的经验。

\eventanchor{JH-0378-BEIFU-ARMY}{太元三年（378年）·谢玄在京口训练北府兵，东晋以侨民和北来军人为核心重组淮南防务}账本记录东晋在太元三年（378年）发生“谢玄在京口训练北府兵，东晋以侨民和北来军人为核心重组淮南防务”。条目把背景概括为前秦进攻梁益和襄阳暴露东晋常备防务不足，谢安需建立能机动应战的部队，并指出北府兵成为淝水之战和其后东晋军政的重要力量，也强化京口军事社会的地位。这个节点进入区域社会时，要经由文书、道路、军队、市场、寺院和家庭等不同中介；因此，政治行动的宣布、地方接收和日常生活的改变不能被视作同一时刻完成。

本条依据《晋书·谢玄传》；《晋书·孝武帝纪》；《资治通鉴·晋纪》，证据提示为“谢玄任职与北府兵形成大体可靠；兵员来源、编制规模和训练时间有争议”。这些材料可支持年序、官方决定、战事或特定地点的线索，却难以直接统计所有人的财富、迁徙、信仰和动机。更稳妥的读法是区分诏令与实施、军报与人口、行政称谓与自我认同，并让地方材料的缺环限制解释的范围。

由此可见，该事件不是孤立的年表点，而是资源、信息与权力关系重新配置的一环。不同地区的官员、社群和家庭可能以不同方式承担征发、保护、迁移或宗教活动；现存来源只能照亮其中一部分，不能替沉默者制造统一的经验。

\eventanchor{JH-0379-XIANGYANG-FALL}{太元四年（379年）·襄阳城陷，守将朱序被俘，前秦控制汉水中游重镇}账本记录前秦与东晋在太元四年（379年）发生“襄阳城陷，守将朱序被俘，前秦控制汉水中游重镇”。条目把背景概括为前秦长期围困并切断外援，东晋难以从江汉方向解除危机，并指出前秦前锋逼近长江流域，淝水决战前的南北力量对比显著恶化。这个节点进入区域社会时，要经由文书、道路、军队、市场、寺院和家庭等不同中介；因此，政治行动的宣布、地方接收和日常生活的改变不能被视作同一时刻完成。

本条依据《晋书·孝武帝纪》；《晋书·朱序传》；《晋书·苻坚载记下》；《资治通鉴·晋纪》，证据提示为“城陷与朱序被俘可靠；守城期限、饥荒情况和投降过程有争议”。这些材料可支持年序、官方决定、战事或特定地点的线索，却难以直接统计所有人的财富、迁徙、信仰和动机。更稳妥的读法是区分诏令与实施、军报与人口、行政称谓与自我认同，并让地方材料的缺环限制解释的范围。

由此可见，该事件不是孤立的年表点，而是资源、信息与权力关系重新配置的一环。不同地区的官员、社群和家庭可能以不同方式承担征发、保护、迁移或宗教活动；现存来源只能照亮其中一部分，不能替沉默者制造统一的经验。

\eventanchor{JH-0382-FU-JIAN-SOUTHERN-INVASION}{太元七年（382年）·苻坚决定大举南伐东晋，调集北方多地兵员与粮运资源}账本记录前秦与东晋在太元七年（382年）发生“苻坚决定大举南伐东晋，调集北方多地兵员与粮运资源”。条目把背景概括为前秦连续兼并北方政权后拥有广阔资源，苻坚高估了对新附地区和将领的控制，并指出全国性征发加重各地负担并暴露统治整合不足，失败风险随之累积。这个节点进入区域社会时，要经由文书、道路、军队、市场、寺院和家庭等不同中介；因此，政治行动的宣布、地方接收和日常生活的改变不能被视作同一时刻完成。

本条依据《晋书·苻坚载记下》；《晋书·孝武帝纪》；《资治通鉴·晋纪》，证据提示为“南伐决策和大规模动员可靠；兵力数字、各族部众的服从程度和王猛遗言叙事有争议”。这些材料可支持年序、官方决定、战事或特定地点的线索，却难以直接统计所有人的财富、迁徙、信仰和动机。更稳妥的读法是区分诏令与实施、军报与人口、行政称谓与自我认同，并让地方材料的缺环限制解释的范围。

由此可见，该事件不是孤立的年表点，而是资源、信息与权力关系重新配置的一环。不同地区的官员、社群和家庭可能以不同方式承担征发、保护、迁移或宗教活动；现存来源只能照亮其中一部分，不能替沉默者制造统一的经验。

\eventanchor{JH-0383-FEI-RIVER}{太元八年（383年）·谢玄、谢石等东晋军在淝水击败前秦南伐军，苻坚北逃}账本记录东晋与前秦在太元八年（383年）发生“谢玄、谢石等东晋军在淝水击败前秦南伐军，苻坚北逃”。条目把背景概括为前秦远征军跨越淮河后补给与指挥困难加剧，东晋北府兵集中于关键渡口，并指出东晋保住江南，前秦权威骤降，北方旧部与新附政权相继脱离。这个节点进入区域社会时，要经由文书、道路、军队、市场、寺院和家庭等不同中介；因此，政治行动的宣布、地方接收和日常生活的改变不能被视作同一时刻完成。

本条依据《晋书·孝武帝纪》；《晋书·谢玄传》；《晋书·苻坚载记下》；《资治通鉴·晋纪》，证据提示为“东晋获胜与前秦军崩溃可靠；兵力、战术口令、风声鹤唳等细节和伤亡数字有争议”。这些材料可支持年序、官方决定、战事或特定地点的线索，却难以直接统计所有人的财富、迁徙、信仰和动机。更稳妥的读法是区分诏令与实施、军报与人口、行政称谓与自我认同，并让地方材料的缺环限制解释的范围。

由此可见，该事件不是孤立的年表点，而是资源、信息与权力关系重新配置的一环。不同地区的官员、社群和家庭可能以不同方式承担征发、保护、迁移或宗教活动；现存来源只能照亮其中一部分，不能替沉默者制造统一的经验。

\eventanchor{JH-0384-MURONG-CHUI-LATER-YAN}{太元九年（384年）·慕容垂在河北起兵，建立后燕并号召慕容旧部反前秦}账本记录后燕在太元九年（384年）发生“慕容垂在河北起兵，建立后燕并号召慕容旧部反前秦”。条目把背景概括为淝水惨败后前秦无力约束前燕旧贵族，慕容垂拥有河北军事与人脉基础，并指出河北出现后燕政权，前秦的北方统治开始碎片化。这个节点进入区域社会时，要经由文书、道路、军队、市场、寺院和家庭等不同中介；因此，政治行动的宣布、地方接收和日常生活的改变不能被视作同一时刻完成。

本条依据《晋书·慕容垂载记》；《资治通鉴·晋纪》；《十六国春秋》辑佚，证据提示为“起兵、建国和慕容旧部响应大体可靠；动员范围、国号使用时点和河北地方态度有争议”。这些材料可支持年序、官方决定、战事或特定地点的线索，却难以直接统计所有人的财富、迁徙、信仰和动机。更稳妥的读法是区分诏令与实施、军报与人口、行政称谓与自我认同，并让地方材料的缺环限制解释的范围。

由此可见，该事件不是孤立的年表点，而是资源、信息与权力关系重新配置的一环。不同地区的官员、社群和家庭可能以不同方式承担征发、保护、迁移或宗教活动；现存来源只能照亮其中一部分，不能替沉默者制造统一的经验。

\eventanchor{JH-0384-YAO-CHANG-LATER-QIN}{太元九年（384年）·姚苌在北地起兵并称万年秦王，后秦政权自前秦军中分出}账本记录后秦在太元九年（384年）发生“姚苌在北地起兵并称万年秦王，后秦政权自前秦军中分出”。条目把背景概括为苻坚南伐失败后无法稳定关中，姚苌与前秦统帅集团关系破裂，并指出后秦占据关中北部并与前秦残余争夺长安，十六国格局更加分散。这个节点进入区域社会时，要经由文书、道路、军队、市场、寺院和家庭等不同中介；因此，政治行动的宣布、地方接收和日常生活的改变不能被视作同一时刻完成。

本条依据《晋书·姚苌载记》；《晋书·苻坚载记下》；《资治通鉴·晋纪》，证据提示为“起兵与称王可靠；羌族部众构成、与苻坚冲突起因和控制地域有争议”。这些材料可支持年序、官方决定、战事或特定地点的线索，却难以直接统计所有人的财富、迁徙、信仰和动机。更稳妥的读法是区分诏令与实施、军报与人口、行政称谓与自我认同，并让地方材料的缺环限制解释的范围。

由此可见，该事件不是孤立的年表点，而是资源、信息与权力关系重新配置的一环。不同地区的官员、社群和家庭可能以不同方式承担征发、保护、迁移或宗教活动；现存来源只能照亮其中一部分，不能替沉默者制造统一的经验。

\eventanchor{JH-0385-FU-JIAN-KILLED}{太元十年（385年）·苻坚在五将山被姚苌部将杀害，前秦失去核心统治者}账本记录前秦与后秦在太元十年（385年）发生“苻坚在五将山被姚苌部将杀害，前秦失去核心统治者”。条目把背景概括为淝水败后苻坚退回关中仍面对叛军与粮饷危机，姚苌掌握独立军队，并指出前秦虽有残余政权但再难恢复统一，后秦取得关中竞争的主动。这个节点进入区域社会时，要经由文书、道路、军队、市场、寺院和家庭等不同中介；因此，政治行动的宣布、地方接收和日常生活的改变不能被视作同一时刻完成。

本条依据《晋书·苻坚载记下》；《晋书·姚苌载记》；《资治通鉴·晋纪》，证据提示为“苻坚被杀和姚苌责任可靠；俘获经过、临终言行和部将执行链条有争议”。这些材料可支持年序、官方决定、战事或特定地点的线索，却难以直接统计所有人的财富、迁徙、信仰和动机。更稳妥的读法是区分诏令与实施、军报与人口、行政称谓与自我认同，并让地方材料的缺环限制解释的范围。

由此可见，该事件不是孤立的年表点，而是资源、信息与权力关系重新配置的一环。不同地区的官员、社群和家庭可能以不同方式承担征发、保护、迁移或宗教活动；现存来源只能照亮其中一部分，不能替沉默者制造统一的经验。

\eventanchor{JH-0386-TUOBA-GUI-RESTORES-WEI}{太元十一年（386年）·拓跋珪在牛川重建拓跋政权并称代王，旋改国号魏}账本记录北魏前身在太元十一年（386年）发生“拓跋珪在牛川重建拓跋政权并称代王，旋改国号魏”。条目把背景概括为前秦崩解造成北部权力真空，拓跋旧部和姻亲网络得以重新集结，并指出北魏成为北方新兴力量，日后将吸纳代北与河北的人口、军队和制度资源。这个节点进入区域社会时，要经由文书、道路、军队、市场、寺院和家庭等不同中介；因此，政治行动的宣布、地方接收和日常生活的改变不能被视作同一时刻完成。

本条依据《魏书·太祖纪》；《魏书·序纪》；《资治通鉴·晋纪》，证据提示为“复国、称王和改魏号大体可靠；早期驻地、部众规模和对前秦遗产的利用有争议”。这些材料可支持年序、官方决定、战事或特定地点的线索，却难以直接统计所有人的财富、迁徙、信仰和动机。更稳妥的读法是区分诏令与实施、军报与人口、行政称谓与自我认同，并让地方材料的缺环限制解释的范围。

由此可见，该事件不是孤立的年表点，而是资源、信息与权力关系重新配置的一环。不同地区的官员、社群和家庭可能以不同方式承担征发、保护、迁移或宗教活动；现存来源只能照亮其中一部分，不能替沉默者制造统一的经验。

\eventanchor{JH-0386-LU-GUANG-LATER-LIANG}{太元十一年（386年）·吕光自西域东归后据凉州建立后凉，河西走廊脱离前秦控制}账本记录后凉在太元十一年（386年）发生“吕光自西域东归后据凉州建立后凉，河西走廊脱离前秦控制”。条目把背景概括为前秦崩解后吕光所部携带独立军事实力返回河西，地方无力抵抗，并指出河西进入多凉政权竞争期，丝路交通与绿洲城市的政治归属频繁变动。这个节点进入区域社会时，要经由文书、道路、军队、市场、寺院和家庭等不同中介；因此，政治行动的宣布、地方接收和日常生活的改变不能被视作同一时刻完成。

本条依据《晋书·吕光载记》；《资治通鉴·晋纪》；《吐鲁番出土文书》，证据提示为“据凉州与后凉建国可靠；西域远征成果、河西城镇归附和国号时点有争议”。这些材料可支持年序、官方决定、战事或特定地点的线索，却难以直接统计所有人的财富、迁徙、信仰和动机。更稳妥的读法是区分诏令与实施、军报与人口、行政称谓与自我认同，并让地方材料的缺环限制解释的范围。

由此可见，该事件不是孤立的年表点，而是资源、信息与权力关系重新配置的一环。不同地区的官员、社群和家庭可能以不同方式承担征发、保护、迁移或宗教活动；现存来源只能照亮其中一部分，不能替沉默者制造统一的经验。

\eventanchor{JH-0386-WESTERN-YAN}{太元十一年（386年）·慕容泓、慕容永等慕容部众在关中、河东活动并建立西燕政权}账本记录西燕在太元十一年（386年）发生“慕容泓、慕容永等慕容部众在关中、河东活动并建立西燕政权”。条目把背景概括为前秦对慕容旧部的控制在淝水后崩溃，部众寻求返回东方故地或夺取新基地，并指出西燕加剧关河地区的军事流动，最终又被后燕整合或击散。这个节点进入区域社会时，要经由文书、道路、军队、市场、寺院和家庭等不同中介；因此，政治行动的宣布、地方接收和日常生活的改变不能被视作同一时刻完成。

本条依据《晋书·慕容泓载记》；《晋书·慕容永载记》；《资治通鉴·晋纪》，证据提示为“西燕形成与慕容永掌权大体可靠；迁徙路线、前后首领更替和部众数量有争议”。这些材料可支持年序、官方决定、战事或特定地点的线索，却难以直接统计所有人的财富、迁徙、信仰和动机。更稳妥的读法是区分诏令与实施、军报与人口、行政称谓与自我认同，并让地方材料的缺环限制解释的范围。

由此可见，该事件不是孤立的年表点，而是资源、信息与权力关系重新配置的一环。不同地区的官员、社群和家庭可能以不同方式承担征发、保护、迁移或宗教活动；现存来源只能照亮其中一部分，不能替沉默者制造统一的经验。

\eventanchor{JH-0392-MURONG-CHUI-DEATH}{太元十七年（392年）·慕容垂去世，慕容宝继位，后燕在北魏与内部矛盾压力下转入不稳}账本记录后燕在太元十七年（392年）发生“慕容垂去世，慕容宝继位，后燕在北魏与内部矛盾压力下转入不稳”。条目把背景概括为慕容垂依靠个人军事威望维系河北，继承人缺乏同等统帅能力，并指出后燕对北魏的优势减弱，河北统治开始出现裂缝。这个节点进入区域社会时，要经由文书、道路、军队、市场、寺院和家庭等不同中介；因此，政治行动的宣布、地方接收和日常生活的改变不能被视作同一时刻完成。

本条依据《晋书·慕容垂载记》；《晋书·慕容宝载记》；《资治通鉴·晋纪》，证据提示为“慕容垂卒与慕容宝继位可靠；遗诏、继承安排和后燕军政资源状况有争议”。这些材料可支持年序、官方决定、战事或特定地点的线索，却难以直接统计所有人的财富、迁徙、信仰和动机。更稳妥的读法是区分诏令与实施、军报与人口、行政称谓与自我认同，并让地方材料的缺环限制解释的范围。

由此可见，该事件不是孤立的年表点，而是资源、信息与权力关系重新配置的一环。不同地区的官员、社群和家庭可能以不同方式承担征发、保护、迁移或宗教活动；现存来源只能照亮其中一部分，不能替沉默者制造统一的经验。

\eventanchor{JH-0393-YAO-XING-ACCESSION}{太元十八年（393年）·姚苌去世，姚兴继位，后秦继续整合关中并调整对前秦残余的战争}账本记录后秦在太元十八年（393年）发生“姚苌去世，姚兴继位，后秦继续整合关中并调整对前秦残余的战争”。条目把背景概括为姚苌以军事联盟建立后秦，继承需要获得羌汉将领和长安官僚承认，并指出姚兴的统治为后秦消灭前秦残余并吸引僧侣、文士创造条件。这个节点进入区域社会时，要经由文书、道路、军队、市场、寺院和家庭等不同中介；因此，政治行动的宣布、地方接收和日常生活的改变不能被视作同一时刻完成。

本条依据《晋书·姚苌载记》；《晋书·姚兴载记》；《资治通鉴·晋纪》，证据提示为“姚苌卒、姚兴继位可靠；继承程序、宗室军权和姚兴早期政策有争议”。这些材料可支持年序、官方决定、战事或特定地点的线索，却难以直接统计所有人的财富、迁徙、信仰和动机。更稳妥的读法是区分诏令与实施、军报与人口、行政称谓与自我认同，并让地方材料的缺环限制解释的范围。

由此可见，该事件不是孤立的年表点，而是资源、信息与权力关系重新配置的一环。不同地区的官员、社群和家庭可能以不同方式承担征发、保护、迁移或宗教活动；现存来源只能照亮其中一部分，不能替沉默者制造统一的经验。

\eventanchor{JH-0394-LATER-QIN-DESTROYS-FORMER-QIN}{太元十九年（394年）·姚兴军击败并杀害苻登，前秦残余政权灭亡}账本记录后秦与前秦残余在太元十九年（394年）发生“姚兴军击败并杀害苻登，前秦残余政权灭亡”。条目把背景概括为苻坚死后前秦数度迁徙且资源匮乏，姚兴继位后集中兵力追击，并指出后秦成为关中主导政权，淝水后残存的前秦帝国结构彻底解体。这个节点进入区域社会时，要经由文书、道路、军队、市场、寺院和家庭等不同中介；因此，政治行动的宣布、地方接收和日常生活的改变不能被视作同一时刻完成。

本条依据《晋书·苻登载记》；《晋书·姚兴载记》；《资治通鉴·晋纪》，证据提示为“苻登败亡与前秦终结可靠；战场、投降者处置和前秦余众去向有争议”。这些材料可支持年序、官方决定、战事或特定地点的线索，却难以直接统计所有人的财富、迁徙、信仰和动机。更稳妥的读法是区分诏令与实施、军报与人口、行政称谓与自我认同，并让地方材料的缺环限制解释的范围。

由此可见，该事件不是孤立的年表点，而是资源、信息与权力关系重新配置的一环。不同地区的官员、社群和家庭可能以不同方式承担征发、保护、迁移或宗教活动；现存来源只能照亮其中一部分，不能替沉默者制造统一的经验。

\eventanchor{JH-0395-CANHEBEI}{太元二十年（395年）·拓跋珪在参合陂大败后燕军，北魏取得对河北竞争的主动}账本记录北魏与后燕在太元二十年（395年）发生“拓跋珪在参合陂大败后燕军，北魏取得对河北竞争的主动”。条目把背景概括为后燕试图以大军压制北魏，行军补给与指挥失误使其暴露于北魏机动部队，并指出后燕精锐受损，北魏由代北政权转为可争夺河北的军事力量。这个节点进入区域社会时，要经由文书、道路、军队、市场、寺院和家庭等不同中介；因此，政治行动的宣布、地方接收和日常生活的改变不能被视作同一时刻完成。

本条依据《魏书·太祖纪》；《晋书·慕容宝载记》；《资治通鉴·晋纪》，证据提示为“战役胜负与后燕军损失严重可靠；伏击地点、兵数和伤亡数字有争议”。这些材料可支持年序、官方决定、战事或特定地点的线索，却难以直接统计所有人的财富、迁徙、信仰和动机。更稳妥的读法是区分诏令与实施、军报与人口、行政称谓与自我认同，并让地方材料的缺环限制解释的范围。

由此可见，该事件不是孤立的年表点，而是资源、信息与权力关系重新配置的一环。不同地区的官员、社群和家庭可能以不同方式承担征发、保护、迁移或宗教活动；现存来源只能照亮其中一部分，不能替沉默者制造统一的经验。

\eventanchor{JH-0396-TUOBA-GUI-EMPEROR}{皇始元年（396年）·拓跋珪称帝，北魏以平城为核心扩张并向后燕发动攻势}账本记录北魏在皇始元年（396年）发生“拓跋珪称帝，北魏以平城为核心扩张并向后燕发动攻势”。条目把背景概括为参合陂胜利提升北魏威望，拓跋珪获得更多部众和城镇归附，并指出北魏以帝号争夺北方正统，后燕的河北控制进一步受压。这个节点进入区域社会时，要经由文书、道路、军队、市场、寺院和家庭等不同中介；因此，政治行动的宣布、地方接收和日常生活的改变不能被视作同一时刻完成。

本条依据《魏书·太祖纪》；《资治通鉴·晋纪》，证据提示为“称帝与北魏扩张可靠；称帝礼仪、都城制度和部落—官僚结构有争议”。这些材料可支持年序、官方决定、战事或特定地点的线索，却难以直接统计所有人的财富、迁徙、信仰和动机。更稳妥的读法是区分诏令与实施、军报与人口、行政称谓与自我认同，并让地方材料的缺环限制解释的范围。

由此可见，该事件不是孤立的年表点，而是资源、信息与权力关系重新配置的一环。不同地区的官员、社群和家庭可能以不同方式承担征发、保护、迁移或宗教活动；现存来源只能照亮其中一部分，不能替沉默者制造统一的经验。

\eventanchor{JH-0397-NORTHERN-WEI-TAKES-ZHONGSHAN}{皇始二年（397年）·北魏攻取中山，后燕政权被迫向东北退却，河北主导权易手}账本记录北魏与后燕在皇始二年（397年）发生“北魏攻取中山，后燕政权被迫向东北退却，河北主导权易手”。条目把背景概括为后燕在参合陂后军力衰弱，慕容宝难以稳定河北城镇，并指出北魏进入中原北部，后燕、南燕和北燕相继分化。这个节点进入区域社会时，要经由文书、道路、军队、市场、寺院和家庭等不同中介；因此，政治行动的宣布、地方接收和日常生活的改变不能被视作同一时刻完成。

本条依据《魏书·太祖纪》；《晋书·慕容宝载记》；《资治通鉴·晋纪》，证据提示为“取中山和后燕退却可靠；城市破坏、投降官民和北魏占领方式有争议”。这些材料可支持年序、官方决定、战事或特定地点的线索，却难以直接统计所有人的财富、迁徙、信仰和动机。更稳妥的读法是区分诏令与实施、军报与人口、行政称谓与自我认同，并让地方材料的缺环限制解释的范围。

由此可见，该事件不是孤立的年表点，而是资源、信息与权力关系重新配置的一环。不同地区的官员、社群和家庭可能以不同方式承担征发、保护、迁移或宗教活动；现存来源只能照亮其中一部分，不能替沉默者制造统一的经验。

\eventanchor{JH-0398-NORTHERN-WEI-MOVES-PINGCHENG}{天兴元年（398年）·拓跋珪迁都平城并重建都城机构，北魏的统治重心由代北部落空间转向城郭行政}账本记录北魏在天兴元年（398年）发生“拓跋珪迁都平城并重建都城机构，北魏的统治重心由代北部落空间转向城郭行政”。条目把背景概括为北魏占有河北后需有更稳定的行政中枢，盛乐旧都难以容纳扩张后的政治需求，并指出平城成为北魏前期军政中心，为后续北方统一和制度建设提供空间条件。这个节点进入区域社会时，要经由文书、道路、军队、市场、寺院和家庭等不同中介；因此，政治行动的宣布、地方接收和日常生活的改变不能被视作同一时刻完成。

本条依据《魏书·太祖纪》；《资治通鉴·晋纪》，证据提示为“迁都与平城经营可靠；宫城布局、迁入人口和制度成熟程度有争议”。这些材料可支持年序、官方决定、战事或特定地点的线索，却难以直接统计所有人的财富、迁徙、信仰和动机。更稳妥的读法是区分诏令与实施、军报与人口、行政称谓与自我认同，并让地方材料的缺环限制解释的范围。

由此可见，该事件不是孤立的年表点，而是资源、信息与权力关系重新配置的一环。不同地区的官员、社群和家庭可能以不同方式承担征发、保护、迁移或宗教活动；现存来源只能照亮其中一部分，不能替沉默者制造统一的经验。

\eventanchor{JH-0399-SUN-EN-REBELLION}{隆安三年（399年）·孙恩在会稽起兵，东晋东南沿海州郡陷入反复征战}账本记录东晋与浙东天师道集团在隆安三年（399年）发生“孙恩在会稽起兵，东晋东南沿海州郡陷入反复征战”。条目把背景概括为东晋财政征发、流民安置和地方豪强矛盾在浙东累积，宗教网络提供动员渠道，并指出东南军政资源被长期牵制，刘裕等将领由平乱获得军事实力。这个节点进入区域社会时，要经由文书、道路、军队、市场、寺院和家庭等不同中介；因此，政治行动的宣布、地方接收和日常生活的改变不能被视作同一时刻完成。

本条依据《晋书·孙恩传》；《晋书·安帝纪》；《资治通鉴·晋纪》，证据提示为“起兵与浙东战乱可靠；天师道组织、信众构成和死亡人数有争议”。这些材料可支持年序、官方决定、战事或特定地点的线索，却难以直接统计所有人的财富、迁徙、信仰和动机。更稳妥的读法是区分诏令与实施、军报与人口、行政称谓与自我认同，并让地方材料的缺环限制解释的范围。

由此可见，该事件不是孤立的年表点，而是资源、信息与权力关系重新配置的一环。不同地区的官员、社群和家庭可能以不同方式承担征发、保护、迁移或宗教活动；现存来源只能照亮其中一部分，不能替沉默者制造统一的经验。

\eventanchor{JH-0401-KUMARAJIVA-CHANGAN}{弘始三年（401年）·鸠摩罗什抵达长安，姚兴支持译经活动，后秦成为重要佛教翻译中心}账本记录后秦在弘始三年（401年）发生“鸠摩罗什抵达长安，姚兴支持译经活动，后秦成为重要佛教翻译中心”。条目把背景概括为后秦控制关中并取得龟兹高僧，姚兴以宗教与文教资源增强长安的吸引力，并指出大量译经和僧侣网络使长安成为跨区域佛教知识传播中心。这个节点进入区域社会时，要经由文书、道路、军队、市场、寺院和家庭等不同中介；因此，政治行动的宣布、地方接收和日常生活的改变不能被视作同一时刻完成。

本条依据《晋书·姚兴载记》；《高僧传·鸠摩罗什传》；《出三藏记集》，证据提示为“抵长安和译场活动可靠；具体入城日期、译场规模和姚兴个人参与程度有争议”。这些材料可支持年序、官方决定、战事或特定地点的线索，却难以直接统计所有人的财富、迁徙、信仰和动机。更稳妥的读法是区分诏令与实施、军报与人口、行政称谓与自我认同，并让地方材料的缺环限制解释的范围。

由此可见，该事件不是孤立的年表点，而是资源、信息与权力关系重新配置的一环。不同地区的官员、社群和家庭可能以不同方式承担征发、保护、迁移或宗教活动；现存来源只能照亮其中一部分，不能替沉默者制造统一的经验。

\subsection{边疆、战争与移动}

\eventanchor{JH-0402-HUAN-XUAN-TAKES-JIANKANG}{元兴元年（402年）·桓玄自荆州东下攻入建康，控制安帝和东晋朝政}账本记录东晋在元兴元年（402年）发生“桓玄自荆州东下攻入建康，控制安帝和东晋朝政”。条目把背景概括为孙恩之乱削弱朝廷，司马道子集团失去军政信誉，桓玄凭荆州兵力进取建康，并指出东晋中枢再度被强藩控制，皇统存续取决于其他军镇的选择。这个节点进入区域社会时，要经由文书、道路、军队、市场、寺院和家庭等不同中介；因此，政治行动的宣布、地方接收和日常生活的改变不能被视作同一时刻完成。

本条依据《晋书·安帝纪》；《晋书·桓玄传》；《资治通鉴·晋纪》；《建康实录》，证据提示为“入建康、控制安帝和桓玄专政可靠；司马道子父子失势过程、城内支持与军队规模有争议”。这些材料可支持年序、官方决定、战事或特定地点的线索，却难以直接统计所有人的财富、迁徙、信仰和动机。更稳妥的读法是区分诏令与实施、军报与人口、行政称谓与自我认同，并让地方材料的缺环限制解释的范围。

由此可见，该事件不是孤立的年表点，而是资源、信息与权力关系重新配置的一环。不同地区的官员、社群和家庭可能以不同方式承担征发、保护、迁移或宗教活动；现存来源只能照亮其中一部分，不能替沉默者制造统一的经验。

\eventanchor{JH-0403-HUAN-XUAN-USURPATION}{元兴二年十二月（403年）·桓玄废安帝自立为楚帝，东晋皇统短暂中断}账本记录桓楚与东晋在元兴二年十二月（403年）发生“桓玄废安帝自立为楚帝，东晋皇统短暂中断”。条目把背景概括为桓玄已控制建康和禁军，试图以帝号将军事优势转化为稳定统治，并指出其僭号激起刘裕等北府将领反弹，桓楚缺乏持久政治联盟。这个节点进入区域社会时，要经由文书、道路、军队、市场、寺院和家庭等不同中介；因此，政治行动的宣布、地方接收和日常生活的改变不能被视作同一时刻完成。

本条依据《晋书·安帝纪》；《晋书·桓玄传》；《资治通鉴·晋纪》，证据提示为“废立与称帝可靠；禅让诏书、朝臣拥戴程度和桓楚制度安排有争议”。这些材料可支持年序、官方决定、战事或特定地点的线索，却难以直接统计所有人的财富、迁徙、信仰和动机。更稳妥的读法是区分诏令与实施、军报与人口、行政称谓与自我认同，并让地方材料的缺环限制解释的范围。

由此可见，该事件不是孤立的年表点，而是资源、信息与权力关系重新配置的一环。不同地区的官员、社群和家庭可能以不同方式承担征发、保护、迁移或宗教活动；现存来源只能照亮其中一部分，不能替沉默者制造统一的经验。

\eventanchor{JH-0404-LIU-YU-RESTORES-JIN}{元兴三年（404年）·刘裕、刘毅等起兵推翻桓楚，迎复安帝，东晋皇统恢复}账本记录东晋在元兴三年（404年）发生“刘裕、刘毅等起兵推翻桓楚，迎复安帝，东晋皇统恢复”。条目把背景概括为桓玄称帝失去北府军和部分士族支持，刘裕利用京口军事网络发动反击，并指出刘裕成为东晋最强军事人物，为其后消灭地方强藩和建宋奠定基础。这个节点进入区域社会时，要经由文书、道路、军队、市场、寺院和家庭等不同中介；因此，政治行动的宣布、地方接收和日常生活的改变不能被视作同一时刻完成。

本条依据《晋书·安帝纪》；《晋书·刘裕传》；《晋书·桓玄传》；《资治通鉴·晋纪》，证据提示为“起兵、桓玄败走和安帝复位可靠；各将领分工、战斗细节和复位礼仪有争议”。这些材料可支持年序、官方决定、战事或特定地点的线索，却难以直接统计所有人的财富、迁徙、信仰和动机。更稳妥的读法是区分诏令与实施、军报与人口、行政称谓与自我认同，并让地方材料的缺环限制解释的范围。

由此可见，该事件不是孤立的年表点，而是资源、信息与权力关系重新配置的一环。不同地区的官员、社群和家庭可能以不同方式承担征发、保护、迁移或宗教活动；现存来源只能照亮其中一部分，不能替沉默者制造统一的经验。

\eventanchor{JH-0405-CUAN-BAOZI-STELE}{义熙元年（405年）·《爨宝子碑》立于宁州，记录爨氏地方官员身份与家族记忆}账本记录东晋宁州地方社会在义熙元年（405年）发生“《爨宝子碑》立于宁州，记录爨氏地方官员身份与家族记忆”。条目把背景概括为东晋中央对西南的控制依赖地方豪族和州郡官员，碑刻成为公开确认身份的媒介，并指出碑文为边远地区的官职、家族与地方秩序提供独立于正史的实物证据。这个节点进入区域社会时，要经由文书、道路、军队、市场、寺院和家庭等不同中介；因此，政治行动的宣布、地方接收和日常生活的改变不能被视作同一时刻完成。

本条依据〈爨宝子碑〉；《晋书·安帝纪》，证据提示为“碑刻纪年、姓名和官衔基本可靠；家世夸饰、官衔性质及其代表的地方控制范围有争议”。这些材料可支持年序、官方决定、战事或特定地点的线索，却难以直接统计所有人的财富、迁徙、信仰和动机。更稳妥的读法是区分诏令与实施、军报与人口、行政称谓与自我认同，并让地方材料的缺环限制解释的范围。

由此可见，该事件不是孤立的年表点，而是资源、信息与权力关系重新配置的一环。不同地区的官员、社群和家庭可能以不同方式承担征发、保护、迁移或宗教活动；现存来源只能照亮其中一部分，不能替沉默者制造统一的经验。

\eventanchor{JH-0406-LIU-YU-DEFEATS-HUAN-REMNANTS}{义熙二年（406年）·刘裕等平定桓玄余部，东晋重新控制荆州与长江中游主要据点}账本记录东晋在义熙二年（406年）发生“刘裕等平定桓玄余部，东晋重新控制荆州与长江中游主要据点”。条目把背景概括为桓玄死后其部将仍据长江中游，东晋复辟必须恢复荆州军政资源，并指出刘裕与同盟将领借平乱分配州镇，东晋后期军阀格局随之重组。这个节点进入区域社会时，要经由文书、道路、军队、市场、寺院和家庭等不同中介；因此，政治行动的宣布、地方接收和日常生活的改变不能被视作同一时刻完成。

本条依据《晋书·刘裕传》；《晋书·刘毅传》；《晋书·安帝纪》；《资治通鉴·晋纪》，证据提示为“桓玄余部败亡和荆州重归东晋大体可靠；余部活动范围、战斗次序和地方响应有争议”。这些材料可支持年序、官方决定、战事或特定地点的线索，却难以直接统计所有人的财富、迁徙、信仰和动机。更稳妥的读法是区分诏令与实施、军报与人口、行政称谓与自我认同，并让地方材料的缺环限制解释的范围。

由此可见，该事件不是孤立的年表点，而是资源、信息与权力关系重新配置的一环。不同地区的官员、社群和家庭可能以不同方式承担征发、保护、迁移或宗教活动；现存来源只能照亮其中一部分，不能替沉默者制造统一的经验。

\eventanchor{JH-0407-NORTHERN-YAN}{义熙三年（407年）·冯跋推翻高云并在龙城建立北燕，后燕残余在东北改组}账本记录北燕在义熙三年（407年）发生“冯跋推翻高云并在龙城建立北燕，后燕残余在东北改组”。条目把背景概括为后燕崩解后东北军政集团围绕龙城重组，高云政权缺乏稳固支持，并指出北燕延续慕容后燕在东北的部分制度和人口基础，并与北魏长期对峙。这个节点进入区域社会时，要经由文书、道路、军队、市场、寺院和家庭等不同中介；因此，政治行动的宣布、地方接收和日常生活的改变不能被视作同一时刻完成。

本条依据《晋书·冯跋载记》；《魏书·冯跋传》；《资治通鉴·晋纪》，证据提示为“冯跋掌权与北燕形成可靠；高云身份、政变参与者和北燕建国时间有争议”。这些材料可支持年序、官方决定、战事或特定地点的线索，却难以直接统计所有人的财富、迁徙、信仰和动机。更稳妥的读法是区分诏令与实施、军报与人口、行政称谓与自我认同，并让地方材料的缺环限制解释的范围。

由此可见，该事件不是孤立的年表点，而是资源、信息与权力关系重新配置的一环。不同地区的官员、社群和家庭可能以不同方式承担征发、保护、迁移或宗教活动；现存来源只能照亮其中一部分，不能替沉默者制造统一的经验。

\eventanchor{JH-0407-HELIAN-BOBO-XIA}{义熙三年（407年）·赫连勃勃在朔方起兵建夏，利用铁弗部及边地军事网络挑战后秦}账本记录夏在义熙三年（407年）发生“赫连勃勃在朔方起兵建夏，利用铁弗部及边地军事网络挑战后秦”。条目把背景概括为后秦对北部边地控制有限，赫连勃勃可凭部族联盟和流动骑兵自立，并指出夏成为关中北部持续威胁，后秦和北魏都须调整北方边防。这个节点进入区域社会时，要经由文书、道路、军队、市场、寺院和家庭等不同中介；因此，政治行动的宣布、地方接收和日常生活的改变不能被视作同一时刻完成。

本条依据《晋书·赫连勃勃载记》；《魏书·铁弗刘虎传》；《资治通鉴·晋纪》，证据提示为“建夏与赫连勃勃掌权可靠；部众来源、国号使用与草原—关中关系有争议”。这些材料可支持年序、官方决定、战事或特定地点的线索，却难以直接统计所有人的财富、迁徙、信仰和动机。更稳妥的读法是区分诏令与实施、军报与人口、行政称谓与自我认同，并让地方材料的缺环限制解释的范围。

由此可见，该事件不是孤立的年表点，而是资源、信息与权力关系重新配置的一环。不同地区的官员、社群和家庭可能以不同方式承担征发、保护、迁移或宗教活动；现存来源只能照亮其中一部分，不能替沉默者制造统一的经验。

\eventanchor{JH-0409-LIU-YU-SOUTHERN-YAN}{义熙五年（409年）·刘裕北伐南燕，东晋军跨淮河进攻广固}账本记录东晋与南燕在义熙五年（409年）发生“刘裕北伐南燕，东晋军跨淮河进攻广固”。条目把背景概括为南燕劫掠东晋边民并控制山东，刘裕希望以北伐巩固个人军政地位，并指出东晋军进入山东，南燕都城遭到长期围攻，北伐声望转化为刘裕的政治资本。这个节点进入区域社会时，要经由文书、道路、军队、市场、寺院和家庭等不同中介；因此，政治行动的宣布、地方接收和日常生活的改变不能被视作同一时刻完成。

本条依据《晋书·刘裕传》；《晋书·慕容超载记》；《资治通鉴·晋纪》，证据提示为“出师、进攻广固和刘裕统帅可靠；军队数目、渡淮路线和南燕守备有争议”。这些材料可支持年序、官方决定、战事或特定地点的线索，却难以直接统计所有人的财富、迁徙、信仰和动机。更稳妥的读法是区分诏令与实施、军报与人口、行政称谓与自我认同，并让地方材料的缺环限制解释的范围。

由此可见，该事件不是孤立的年表点，而是资源、信息与权力关系重新配置的一环。不同地区的官员、社群和家庭可能以不同方式承担征发、保护、迁移或宗教活动；现存来源只能照亮其中一部分，不能替沉默者制造统一的经验。

\eventanchor{JH-0410-SOUTHERN-YAN-FALL}{义熙六年（410年）·刘裕攻陷广固，俘杀慕容超，南燕灭亡}账本记录东晋与南燕在义熙六年（410年）发生“刘裕攻陷广固，俘杀慕容超，南燕灭亡”。条目把背景概括为广固久围后粮食与援军均告匮乏，南燕无法摆脱东晋军的封锁，并指出东晋短暂取得山东部分地区，刘裕在南方朝廷中的威望进一步上升。这个节点进入区域社会时，要经由文书、道路、军队、市场、寺院和家庭等不同中介；因此，政治行动的宣布、地方接收和日常生活的改变不能被视作同一时刻完成。

本条依据《晋书·刘裕传》；《晋书·慕容超载记》；《资治通鉴·晋纪》，证据提示为“广固陷落、慕容超被俘和南燕灭亡可靠；围城损失、战后屠掠和山东接管程度有争议”。这些材料可支持年序、官方决定、战事或特定地点的线索，却难以直接统计所有人的财富、迁徙、信仰和动机。更稳妥的读法是区分诏令与实施、军报与人口、行政称谓与自我认同，并让地方材料的缺环限制解释的范围。

由此可见，该事件不是孤立的年表点，而是资源、信息与权力关系重新配置的一环。不同地区的官员、社群和家庭可能以不同方式承担征发、保护、迁移或宗教活动；现存来源只能照亮其中一部分，不能替沉默者制造统一的经验。

\eventanchor{JH-0412-LIU-YU-DEFEATS-LIU-YI}{义熙八年（412年）·刘裕击败并杀害据荆州的刘毅，将长江中游强镇纳入自身控制}账本记录东晋在义熙八年（412年）发生“刘裕击败并杀害据荆州的刘毅，将长江中游强镇纳入自身控制”。条目把背景概括为刘裕与刘毅同为反桓玄功臣，二人对荆州资源和朝廷主导权的竞争不可调和，并指出刘裕消除主要军事对手，得以更直接调度长江中游兵粮进行北伐。这个节点进入区域社会时，要经由文书、道路、军队、市场、寺院和家庭等不同中介；因此，政治行动的宣布、地方接收和日常生活的改变不能被视作同一时刻完成。

本条依据《晋书·刘裕传》；《晋书·刘毅传》；《资治通鉴·晋纪》，证据提示为“刘毅败亡与刘裕接收荆州可靠；冲突起因、各州将领选择和清算范围有争议”。这些材料可支持年序、官方决定、战事或特定地点的线索，却难以直接统计所有人的财富、迁徙、信仰和动机。更稳妥的读法是区分诏令与实施、军报与人口、行政称谓与自我认同，并让地方材料的缺环限制解释的范围。

由此可见，该事件不是孤立的年表点，而是资源、信息与权力关系重新配置的一环。不同地区的官员、社群和家庭可能以不同方式承担征发、保护、迁移或宗教活动；现存来源只能照亮其中一部分，不能替沉默者制造统一的经验。

\eventanchor{JH-0413-QIAO-ZONG-DEFEATED}{义熙九年（413年）·朱龄石等东晋军攻灭谯蜀，成都重新纳入东晋治下}账本记录东晋与谯蜀在义熙九年（413年）发生“朱龄石等东晋军攻灭谯蜀，成都重新纳入东晋治下”。条目把背景概括为刘裕削弱刘毅后可调动荆州兵力入蜀，谯蜀内部与后秦关系亦趋脆弱，并指出东晋恢复对成都的控制，长江上游资源再次进入刘裕的军事体系。这个节点进入区域社会时，要经由文书、道路、军队、市场、寺院和家庭等不同中介；因此，政治行动的宣布、地方接收和日常生活的改变不能被视作同一时刻完成。

本条依据《晋书·朱龄石传》；《晋书·谯纵载记》；《资治通鉴·晋纪》，证据提示为“谯蜀灭亡和成都收复可靠；入蜀路线、谯纵死因及地方接收过程有争议”。这些材料可支持年序、官方决定、战事或特定地点的线索，却难以直接统计所有人的财富、迁徙、信仰和动机。更稳妥的读法是区分诏令与实施、军报与人口、行政称谓与自我认同，并让地方材料的缺环限制解释的范围。

由此可见，该事件不是孤立的年表点，而是资源、信息与权力关系重新配置的一环。不同地区的官员、社群和家庭可能以不同方式承担征发、保护、迁移或宗教活动；现存来源只能照亮其中一部分，不能替沉默者制造统一的经验。

\eventanchor{JH-0413-LU-XUN-SUPPRESSED}{义熙九年（413年）·卢循在广州一带败亡，孙恩—卢循集团的长期反乱结束}账本记录东晋与岭南地方集团在义熙九年（413年）发生“卢循在广州一带败亡，孙恩—卢循集团的长期反乱结束”。条目把背景概括为卢循北上失败后退入岭南，东晋以水军和地方将领持续追击，并指出东晋解除东南海路与岭南的重大威胁，刘裕可集中资源处理北方和朝廷政局。这个节点进入区域社会时，要经由文书、道路、军队、市场、寺院和家庭等不同中介；因此，政治行动的宣布、地方接收和日常生活的改变不能被视作同一时刻完成。

本条依据《晋书·卢循传》；《晋书·刘裕传》；《资治通鉴·晋纪》，证据提示为“卢循败亡与叛乱终结可靠；自尽地点、部众去向和岭南地方社会参与有争议”。这些材料可支持年序、官方决定、战事或特定地点的线索，却难以直接统计所有人的财富、迁徙、信仰和动机。更稳妥的读法是区分诏令与实施、军报与人口、行政称谓与自我认同，并让地方材料的缺环限制解释的范围。

由此可见，该事件不是孤立的年表点，而是资源、信息与权力关系重新配置的一环。不同地区的官员、社群和家庭可能以不同方式承担征发、保护、迁移或宗教活动；现存来源只能照亮其中一部分，不能替沉默者制造统一的经验。

\eventanchor{JH-0414-GWANGGAETO-STELE}{义熙十年（414年）·高句丽为好太王立碑，铭刻其征伐、城郭和王权记忆}账本记录高句丽与东北诸政权在义熙十年（414年）发生“高句丽为好太王立碑，铭刻其征伐、城郭和王权记忆”。条目把背景概括为高句丽在东北扩张后需要以碑刻公开建构王权与征服记忆，并指出碑文提供与中原正史不同的东北区域视角，但其政治修辞须与其他材料互证。这个节点进入区域社会时，要经由文书、道路、军队、市场、寺院和家庭等不同中介；因此，政治行动的宣布、地方接收和日常生活的改变不能被视作同一时刻完成。

本条依据〈好太王碑〉；《三国史记·广开土王本纪》；《魏书·高句丽传》，证据提示为“立碑年代与主体内容大体可靠；残损字句、倭相关段落和战争叙事的夸张程度有争议”。这些材料可支持年序、官方决定、战事或特定地点的线索，却难以直接统计所有人的财富、迁徙、信仰和动机。更稳妥的读法是区分诏令与实施、军报与人口、行政称谓与自我认同，并让地方材料的缺环限制解释的范围。

由此可见，该事件不是孤立的年表点，而是资源、信息与权力关系重新配置的一环。不同地区的官员、社群和家庭可能以不同方式承担征发、保护、迁移或宗教活动；现存来源只能照亮其中一部分，不能替沉默者制造统一的经验。

\eventanchor{JH-0414-WESTERN-QIN-DESTROYS-SOUTHERN-LIANG}{义熙十年（414年）·乞伏炽磐攻灭南凉，西秦在河湟地区扩展控制}账本记录西秦与南凉在义熙十年（414年）发生“乞伏炽磐攻灭南凉，西秦在河湟地区扩展控制”。条目把背景概括为河西诸凉彼此争战，南凉资源与联盟逐渐耗尽，西秦趁机东进，并指出河湟力量格局改变，西秦成为后秦、夏与北凉之间的重要政权。这个节点进入区域社会时，要经由文书、道路、军队、市场、寺院和家庭等不同中介；因此，政治行动的宣布、地方接收和日常生活的改变不能被视作同一时刻完成。

本条依据《晋书·乞伏炽磐载记》；《晋书·秃发傉檀载记》；《资治通鉴·晋纪》，证据提示为“南凉灭亡与西秦扩张可靠；战役路线、秃发氏部众去向和河湟城镇控制有争议”。这些材料可支持年序、官方决定、战事或特定地点的线索，却难以直接统计所有人的财富、迁徙、信仰和动机。更稳妥的读法是区分诏令与实施、军报与人口、行政称谓与自我认同，并让地方材料的缺环限制解释的范围。

由此可见，该事件不是孤立的年表点，而是资源、信息与权力关系重新配置的一环。不同地区的官员、社群和家庭可能以不同方式承担征发、保护、迁移或宗教活动；现存来源只能照亮其中一部分，不能替沉默者制造统一的经验。

\eventanchor{JH-0416-LIU-YU-NORTHERN-EXPEDITION}{义熙十二年（416年）·刘裕自建康发动大规模北伐，水陆并进攻向后秦控制的河南与关中}账本记录东晋与后秦在义熙十二年（416年）发生“刘裕自建康发动大规模北伐，水陆并进攻向后秦控制的河南与关中”。条目把背景概括为刘裕已控制荆州、益州等资源，后秦又因继承与夏国压力衰弱，并指出东晋军越过淮河与汉水，后秦的河南、关中防线同时承压。这个节点进入区域社会时，要经由文书、道路、军队、市场、寺院和家庭等不同中介；因此，政治行动的宣布、地方接收和日常生活的改变不能被视作同一时刻完成。

本条依据《晋书·刘裕传》；《晋书·安帝纪》；《晋书·姚泓载记》；《资治通鉴·晋纪》，证据提示为“出师、主要将领和多路进军可靠；兵员数字、补给规模和北伐目标层级有争议”。这些材料可支持年序、官方决定、战事或特定地点的线索，却难以直接统计所有人的财富、迁徙、信仰和动机。更稳妥的读法是区分诏令与实施、军报与人口、行政称谓与自我认同，并让地方材料的缺环限制解释的范围。

由此可见，该事件不是孤立的年表点，而是资源、信息与权力关系重新配置的一环。不同地区的官员、社群和家庭可能以不同方式承担征发、保护、迁移或宗教活动；现存来源只能照亮其中一部分，不能替沉默者制造统一的经验。

\eventanchor{JH-0417-LIU-YU-TAKES-LUOYANG}{义熙十三年（417年）·刘裕军攻取洛阳，后秦在河南的防线崩解}账本记录东晋与后秦在义熙十三年（417年）发生“刘裕军攻取洛阳，后秦在河南的防线崩解”。条目把背景概括为东晋军多路北进并切断后秦河南援军，洛阳守备难以久持，并指出刘裕获得进攻关中的前进基地，也以收复旧都加强其政治威望。这个节点进入区域社会时，要经由文书、道路、军队、市场、寺院和家庭等不同中介；因此，政治行动的宣布、地方接收和日常生活的改变不能被视作同一时刻完成。

本条依据《晋书·刘裕传》；《晋书·姚泓载记》；《资治通鉴·晋纪》，证据提示为“取洛阳和后秦退守关中可靠；城内交接、东晋驻军规模和地方民众态度有争议”。这些材料可支持年序、官方决定、战事或特定地点的线索，却难以直接统计所有人的财富、迁徙、信仰和动机。更稳妥的读法是区分诏令与实施、军报与人口、行政称谓与自我认同，并让地方材料的缺环限制解释的范围。

由此可见，该事件不是孤立的年表点，而是资源、信息与权力关系重新配置的一环。不同地区的官员、社群和家庭可能以不同方式承担征发、保护、迁移或宗教活动；现存来源只能照亮其中一部分，不能替沉默者制造统一的经验。

\eventanchor{JH-0417-LIU-YU-TAKES-CHANGAN}{义熙十三年（417年）·东晋军攻入长安，姚泓降服，后秦灭亡}账本记录东晋与后秦在义熙十三年（417年）发生“东晋军攻入长安，姚泓降服，后秦灭亡”。条目把背景概括为后秦遭东晋与夏国两面压力，姚泓无法整合关中将领和城防资源，并指出东晋短暂恢复长安与关中，但远离江南中枢的占领难以长期维持。这个节点进入区域社会时，要经由文书、道路、军队、市场、寺院和家庭等不同中介；因此，政治行动的宣布、地方接收和日常生活的改变不能被视作同一时刻完成。

本条依据《晋书·刘裕传》；《晋书·姚泓载记》；《资治通鉴·晋纪》，证据提示为“长安陷落、姚泓降服和后秦终结可靠；攻城过程、战后处置和东晋对关中的控制深度有争议”。这些材料可支持年序、官方决定、战事或特定地点的线索，却难以直接统计所有人的财富、迁徙、信仰和动机。更稳妥的读法是区分诏令与实施、军报与人口、行政称谓与自我认同，并让地方材料的缺环限制解释的范围。

由此可见，该事件不是孤立的年表点，而是资源、信息与权力关系重新配置的一环。不同地区的官员、社群和家庭可能以不同方式承担征发、保护、迁移或宗教活动；现存来源只能照亮其中一部分，不能替沉默者制造统一的经验。

\eventanchor{JH-0418-XIA-TAKES-CHANGAN}{义熙十四年（418年）·赫连勃勃趁刘裕南返攻取长安，东晋失去关中据点}账本记录夏与东晋关中守军在义熙十四年（418年）发生“赫连勃勃趁刘裕南返攻取长安，东晋失去关中据点”。条目把背景概括为刘裕须返回建康处理皇位与朝政，留守军队兵少粮乏，夏军则熟悉关中北部，并指出东晋的关中占领迅速结束，夏国获得长安并成为西北强权。这个节点进入区域社会时，要经由文书、道路、军队、市场、寺院和家庭等不同中介；因此，政治行动的宣布、地方接收和日常生活的改变不能被视作同一时刻完成。

本条依据《晋书·赫连勃勃载记》；《晋书·刘裕传》；《资治通鉴·晋纪》，证据提示为“夏军取长安和东晋守军撤离可靠；刘裕南返动机、守将责任和城中交接有争议”。这些材料可支持年序、官方决定、战事或特定地点的线索，却难以直接统计所有人的财富、迁徙、信仰和动机。更稳妥的读法是区分诏令与实施、军报与人口、行政称谓与自我认同，并让地方材料的缺环限制解释的范围。

由此可见，该事件不是孤立的年表点，而是资源、信息与权力关系重新配置的一环。不同地区的官员、社群和家庭可能以不同方式承担征发、保护、迁移或宗教活动；现存来源只能照亮其中一部分，不能替沉默者制造统一的经验。

\eventanchor{JH-0420-LIU-YU-FOUNDS-SONG}{永初元年六月（420年）·刘裕受禅建立宋，东晋皇统终结，南朝刘宋开始}账本记录刘宋与东晋在永初元年六月（420年）发生“刘裕受禅建立宋，东晋皇统终结，南朝刘宋开始”。条目把背景概括为刘裕在北伐与平定强藩后掌握朝廷和主要军队，东晋恭帝缺乏独立政治基础，并指出江南皇统转入刘宋，东晋时期的侨姓与军镇结构在新王朝中重新调整。这个节点进入区域社会时，要经由文书、道路、军队、市场、寺院和家庭等不同中介；因此，政治行动的宣布、地方接收和日常生活的改变不能被视作同一时刻完成。

本条依据《宋书·武帝纪》；《晋书·恭帝纪》；《资治通鉴·晋纪》，证据提示为“受禅、改元和东晋终结可靠；禅让文书、恭帝处境和朝臣参与程度有争议”。这些材料可支持年序、官方决定、战事或特定地点的线索，却难以直接统计所有人的财富、迁徙、信仰和动机。更稳妥的读法是区分诏令与实施、军报与人口、行政称谓与自我认同，并让地方材料的缺环限制解释的范围。

由此可见，该事件不是孤立的年表点，而是资源、信息与权力关系重新配置的一环。不同地区的官员、社群和家庭可能以不同方式承担征发、保护、迁移或宗教活动；现存来源只能照亮其中一部分，不能替沉默者制造统一的经验。

\eventanchor{JH-0420-BINGLING-CAVE-169}{建弘元年（420年）·炳灵寺第169窟题记留下西秦年号与佛教造像活动的纪年证据}账本记录西秦在建弘元年（420年）发生“炳灵寺第169窟题记留下西秦年号与佛教造像活动的纪年证据”。条目把背景概括为河湟政权竞争中，地方官民与僧团通过石窟造像建立功德和公共记忆，并指出题记为西秦年号、宗教网络和区域物质文化提供可定位的考古锚点。这个节点进入区域社会时，要经由文书、道路、军队、市场、寺院和家庭等不同中介；因此，政治行动的宣布、地方接收和日常生活的改变不能被视作同一时刻完成。

本条依据炳灵寺第169窟西秦建弘元年题记；《晋书·乞伏炽磐载记》；炳灵寺石窟考古报告，证据提示为“题记年号和洞窟活动可靠；营建主体、造像阶段与国家赞助程度有争议”。这些材料可支持年序、官方决定、战事或特定地点的线索，却难以直接统计所有人的财富、迁徙、信仰和动机。更稳妥的读法是区分诏令与实施、军报与人口、行政称谓与自我认同，并让地方材料的缺环限制解释的范围。

由此可见，该事件不是孤立的年表点，而是资源、信息与权力关系重新配置的一环。不同地区的官员、社群和家庭可能以不同方式承担征发、保护、迁移或宗教活动；现存来源只能照亮其中一部分，不能替沉默者制造统一的经验。

\eventanchor{JH-0422-MINGYUAN-EMPEROR-DEATH}{泰常七年（422年）·明元帝拓跋嗣去世，太武帝拓跋焘继位，北魏军政权力完成交接}账本记录北魏在泰常七年（422年）发生“明元帝拓跋嗣去世，太武帝拓跋焘继位，北魏军政权力完成交接”。条目把背景概括为北魏已控制北方大片地区，皇位继承需平衡宗室、代北贵族与汉人官僚，并指出太武帝继位开启北魏进一步统一北方和强化军政的阶段。这个节点进入区域社会时，要经由文书、道路、军队、市场、寺院和家庭等不同中介；因此，政治行动的宣布、地方接收和日常生活的改变不能被视作同一时刻完成。

本条依据《魏书·明元帝纪》；《魏书·世祖纪》；《资治通鉴·宋纪》，证据提示为“明元帝卒与太武帝继位可靠；遗诏、辅政大臣角色和新帝早期自主性有争议”。这些材料可支持年序、官方决定、战事或特定地点的线索，却难以直接统计所有人的财富、迁徙、信仰和动机。更稳妥的读法是区分诏令与实施、军报与人口、行政称谓与自我认同，并让地方材料的缺环限制解释的范围。

由此可见，该事件不是孤立的年表点，而是资源、信息与权力关系重新配置的一环。不同地区的官员、社群和家庭可能以不同方式承担征发、保护、迁移或宗教活动；现存来源只能照亮其中一部分，不能替沉默者制造统一的经验。

\eventanchor{JH-0424-TAIWU-CONSOLIDATES}{始光元年（424年）·太武帝亲政后处置权臣与宗室竞争，重新集中北魏军政指挥}账本记录北魏在始光元年（424年）发生“太武帝亲政后处置权臣与宗室竞争，重新集中北魏军政指挥”。条目把背景概括为幼帝继位时辅政集团掌权，太武帝需建立对禁军和外镇的直接控制，并指出北魏得以更集中地发动对夏、北燕和北凉的战争。这个节点进入区域社会时，要经由文书、道路、军队、市场、寺院和家庭等不同中介；因此，政治行动的宣布、地方接收和日常生活的改变不能被视作同一时刻完成。

本条依据《魏书·世祖纪》；《魏书·崔浩传》；《资治通鉴·宋纪》，证据提示为“太武帝亲政与政治整顿大体可靠；权臣罪状、清洗范围和决策先后有争议”。这些材料可支持年序、官方决定、战事或特定地点的线索，却难以直接统计所有人的财富、迁徙、信仰和动机。更稳妥的读法是区分诏令与实施、军报与人口、行政称谓与自我认同，并让地方材料的缺环限制解释的范围。

由此可见，该事件不是孤立的年表点，而是资源、信息与权力关系重新配置的一环。不同地区的官员、社群和家庭可能以不同方式承担征发、保护、迁移或宗教活动；现存来源只能照亮其中一部分，不能替沉默者制造统一的经验。

\eventanchor{JH-0427-NORTHERN-WEI-TAKES-TONGWAN}{始光四年（427年）·太武帝攻取夏国都城统万城，夏政权的核心基地失守}账本记录北魏与夏在始光四年（427年）发生“太武帝攻取夏国都城统万城，夏政权的核心基地失守”。条目把背景概括为夏国在赫连勃勃死后继承不稳，北魏已完成对关中北部的战略包围，并指出夏失去都城与军需中心，北魏控制陕北通道并加速统一北方。这个节点进入区域社会时，要经由文书、道路、军队、市场、寺院和家庭等不同中介；因此，政治行动的宣布、地方接收和日常生活的改变不能被视作同一时刻完成。

本条依据《魏书·世祖纪》；《晋书·赫连勃勃载记》；《资治通鉴·宋纪》，证据提示为“统万城陷落与北魏获胜可靠；城防破坏、赫连昌去向和俘虏处置有争议”。这些材料可支持年序、官方决定、战事或特定地点的线索，却难以直接统计所有人的财富、迁徙、信仰和动机。更稳妥的读法是区分诏令与实施、军报与人口、行政称谓与自我认同，并让地方材料的缺环限制解释的范围。

由此可见，该事件不是孤立的年表点，而是资源、信息与权力关系重新配置的一环。不同地区的官员、社群和家庭可能以不同方式承担征发、保护、迁移或宗教活动；现存来源只能照亮其中一部分，不能替沉默者制造统一的经验。

\eventanchor{JH-0430-LIU-SONG-NORTHERN-EXPEDITION}{元嘉七年（430年）·刘宋北伐一度收复黄河以南若干城镇，旋因北魏反击而难以久守}账本记录刘宋与北魏在元嘉七年（430年）发生“刘宋北伐一度收复黄河以南若干城镇，旋因北魏反击而难以久守”。条目把背景概括为刘宋希望恢复刘裕北伐遗产，北魏则以黄河为战略屏障部署反击，并指出南朝难以稳固经营河南，南北对峙的地理边界逐渐清晰。这个节点进入区域社会时，要经由文书、道路、军队、市场、寺院和家庭等不同中介；因此，政治行动的宣布、地方接收和日常生活的改变不能被视作同一时刻完成。

本条依据《宋书·文帝纪》；《宋书·到彦之传》；《魏书·世祖纪》；《资治通鉴·宋纪》，证据提示为“北伐、短暂收复与后撤大体可靠；收复地域、军力规模和撤军责任有争议”。这些材料可支持年序、官方决定、战事或特定地点的线索，却难以直接统计所有人的财富、迁徙、信仰和动机。更稳妥的读法是区分诏令与实施、军报与人口、行政称谓与自我认同，并让地方材料的缺环限制解释的范围。

由此可见，该事件不是孤立的年表点，而是资源、信息与权力关系重新配置的一环。不同地区的官员、社群和家庭可能以不同方式承担征发、保护、迁移或宗教活动；现存来源只能照亮其中一部分，不能替沉默者制造统一的经验。

\eventanchor{JH-0431-XIA-DESTROYED}{神麚四年（431年）·赫连定败亡，夏国残余政权终结，北魏进一步控制关陇北部}账本记录北魏与夏残余在神麚四年（431年）发生“赫连定败亡，夏国残余政权终结，北魏进一步控制关陇北部”。条目把背景概括为统万城失守后夏军西迁，失去稳定粮源和城郭基地，且受北魏与吐谷浑夹击，并指出北魏消除关中北部主要对手，得以把兵力转向北燕和河西。这个节点进入区域社会时，要经由文书、道路、军队、市场、寺院和家庭等不同中介；因此，政治行动的宣布、地方接收和日常生活的改变不能被视作同一时刻完成。

本条依据《魏书·世祖纪》；《晋书·赫连勃勃载记》；《资治通鉴·宋纪》，证据提示为“夏残余败亡与政权终结大体可靠；赫连定被擒经过、吐谷浑角色和终结年份界定有争议”。这些材料可支持年序、官方决定、战事或特定地点的线索，却难以直接统计所有人的财富、迁徙、信仰和动机。更稳妥的读法是区分诏令与实施、军报与人口、行政称谓与自我认同，并让地方材料的缺环限制解释的范围。

由此可见，该事件不是孤立的年表点，而是资源、信息与权力关系重新配置的一环。不同地区的官员、社群和家庭可能以不同方式承担征发、保护、迁移或宗教活动；现存来源只能照亮其中一部分，不能替沉默者制造统一的经验。

\eventanchor{JH-0433-HELIAN-DING-CAPTURED}{延和二年（433年）·吐谷浑将赫连定送交北魏，夏国最后君主的流亡政治结束}账本记录北魏、夏残余与吐谷浑在延和二年（433年）发生“吐谷浑将赫连定送交北魏，夏国最后君主的流亡政治结束”。条目把背景概括为夏残余脱离统万城后依赖西北流动路线，难以抵抗北魏与吐谷浑协作，并指出北魏以收押赫连定宣告夏的彻底终结，并强化对西北诸势力的威慑。这个节点进入区域社会时，要经由文书、道路、军队、市场、寺院和家庭等不同中介；因此，政治行动的宣布、地方接收和日常生活的改变不能被视作同一时刻完成。

本条依据《魏书·世祖纪》；《资治通鉴·宋纪》；《北史·吐谷浑传》，证据提示为“赫连定被送北魏可靠；被擒地点、吐谷浑与北魏交涉及夏残部去向有争议”。这些材料可支持年序、官方决定、战事或特定地点的线索，却难以直接统计所有人的财富、迁徙、信仰和动机。更稳妥的读法是区分诏令与实施、军报与人口、行政称谓与自我认同，并让地方材料的缺环限制解释的范围。

由此可见，该事件不是孤立的年表点，而是资源、信息与权力关系重新配置的一环。不同地区的官员、社群和家庭可能以不同方式承担征发、保护、迁移或宗教活动；现存来源只能照亮其中一部分，不能替沉默者制造统一的经验。

\eventanchor{JH-0436-NORTHERN-WEI-DESTROYS-NORTHERN-YAN}{太延二年（436年）·北魏攻灭北燕，冯弘出奔高句丽，龙城政权结束}账本记录北魏与北燕在太延二年（436年）发生“北魏攻灭北燕，冯弘出奔高句丽，龙城政权结束”。条目把背景概括为北燕长期受北魏军事压力且难获南方援助，高句丽成为其最后外部选择，并指出北魏消除东北腹地对手，但与高句丽的边境关系更趋紧张。这个节点进入区域社会时，要经由文书、道路、军队、市场、寺院和家庭等不同中介；因此，政治行动的宣布、地方接收和日常生活的改变不能被视作同一时刻完成。

本条依据《魏书·世祖纪》；《魏书·冯弘传》；《资治通鉴·宋纪》，证据提示为“北燕灭亡与冯弘出奔可靠；攻城过程、高句丽接纳原因和北燕遗民去向有争议”。这些材料可支持年序、官方决定、战事或特定地点的线索，却难以直接统计所有人的财富、迁徙、信仰和动机。更稳妥的读法是区分诏令与实施、军报与人口、行政称谓与自我认同，并让地方材料的缺环限制解释的范围。

由此可见，该事件不是孤立的年表点，而是资源、信息与权力关系重新配置的一环。不同地区的官员、社群和家庭可能以不同方式承担征发、保护、迁移或宗教活动；现存来源只能照亮其中一部分，不能替沉默者制造统一的经验。

\eventanchor{JH-0439-NORTHERN-WEI-DESTROYS-NORTHERN-LIANG}{太延五年（439年）·北魏攻取姑臧并灭北凉，北方主要割据政权被纳入北魏统治}账本记录北魏与北凉在太延五年（439年）发生“北魏攻取姑臧并灭北凉，北方主要割据政权被纳入北魏统治”。条目把背景概括为北魏先后消灭夏与北燕后可集中西进，北凉又因继承和外援不足难以持久，并指出北魏完成对北方主要政权的军事整合，但河西、西域与地方社会的实际纳入仍需长期经营。这个节点进入区域社会时，要经由文书、道路、军队、市场、寺院和家庭等不同中介；因此，政治行动的宣布、地方接收和日常生活的改变不能被视作同一时刻完成。

本条依据《魏书·世祖纪》；《魏书·沮渠蒙逊传》；《资治通鉴·宋纪》，证据提示为“北凉灭亡、姑臧失守和北方统一结果可靠；沮渠氏出奔、河西城镇接收和统一范围有争议”。这些材料可支持年序、官方决定、战事或特定地点的线索，却难以直接统计所有人的财富、迁徙、信仰和动机。更稳妥的读法是区分诏令与实施、军报与人口、行政称谓与自我认同，并让地方材料的缺环限制解释的范围。

由此可见，该事件不是孤立的年表点，而是资源、信息与权力关系重新配置的一环。不同地区的官员、社群和家庭可能以不同方式承担征发、保护、迁移或宗教活动；现存来源只能照亮其中一部分，不能替沉默者制造统一的经验。

\subsection{社会关系与宗教空间}

\eventanchor{JH-LEDGER-019-1}{永嘉元年（307年）·“司马睿出镇建业，王导等协助联络江东士族并重整东南军政”的前期动员与资源准备}账本记录西晋江东在永嘉元年（307年）发生““司马睿出镇建业，王导等协助联络江东士族并重整东南军政”的前期动员与资源准备”。条目把背景概括为当时的权力、资源与信息约束推动相关过程展开，并指出这一过程为“司马睿出镇建业，王导等协助联络江东士族并重整东南军政”提供可观察的条件。这个节点进入区域社会时，要经由文书、道路、军队、市场、寺院和家庭等不同中介；因此，政治行动的宣布、地方接收和日常生活的改变不能被视作同一时刻完成。

本条依据《晋书·元帝纪》；《晋书·王导传》；《建康实录》，证据提示为“核心事项已有既有账本来源；本分解节点的参与者、执行范围和时间先后仍须逐案核验”。这些材料可支持年序、官方决定、战事或特定地点的线索，却难以直接统计所有人的财富、迁徙、信仰和动机。更稳妥的读法是区分诏令与实施、军报与人口、行政称谓与自我认同，并让地方材料的缺环限制解释的范围。

由此可见，该事件不是孤立的年表点，而是资源、信息与权力关系重新配置的一环。不同地区的官员、社群和家庭可能以不同方式承担征发、保护、迁移或宗教活动；现存来源只能照亮其中一部分，不能替沉默者制造统一的经验。

\eventanchor{JH-LEDGER-019-2}{永嘉元年（307年）·“司马睿出镇建业，王导等协助联络江东士族并重整东南军政”的命令、联盟或地方响应}账本记录西晋江东在永嘉元年（307年）发生““司马睿出镇建业，王导等协助联络江东士族并重整东南军政”的命令、联盟或地方响应”。条目把背景概括为当时的权力、资源与信息约束推动相关过程展开，并指出这一过程为“司马睿出镇建业，王导等协助联络江东士族并重整东南军政”提供可观察的条件。这个节点进入区域社会时，要经由文书、道路、军队、市场、寺院和家庭等不同中介；因此，政治行动的宣布、地方接收和日常生活的改变不能被视作同一时刻完成。

本条依据《晋书·元帝纪》；《晋书·王导传》；《建康实录》，证据提示为“核心事项已有既有账本来源；本分解节点的参与者、执行范围和时间先后仍须逐案核验”。这些材料可支持年序、官方决定、战事或特定地点的线索，却难以直接统计所有人的财富、迁徙、信仰和动机。更稳妥的读法是区分诏令与实施、军报与人口、行政称谓与自我认同，并让地方材料的缺环限制解释的范围。

由此可见，该事件不是孤立的年表点，而是资源、信息与权力关系重新配置的一环。不同地区的官员、社群和家庭可能以不同方式承担征发、保护、迁移或宗教活动；现存来源只能照亮其中一部分，不能替沉默者制造统一的经验。

\eventanchor{JH-LEDGER-019-3}{永嘉元年（307年）·司马睿出镇建业，王导等协助联络江东士族并重整东南军政}账本记录西晋江东在永嘉元年（307年）发生“司马睿出镇建业，王导等协助联络江东士族并重整东南军政”。条目把背景概括为“司马睿出镇建业，王导等协助联络江东士族并重整东南军政”发生的政治与区域条件共同作用，并指出该节点改变后续资源、联盟或制度处置的条件。这个节点进入区域社会时，要经由文书、道路、军队、市场、寺院和家庭等不同中介；因此，政治行动的宣布、地方接收和日常生活的改变不能被视作同一时刻完成。

本条依据《晋书·元帝纪》；《晋书·王导传》；《建康实录》，证据提示为“核心事项已有既有账本来源；本分解节点的参与者、执行范围和时间先后仍须逐案核验”。这些材料可支持年序、官方决定、战事或特定地点的线索，却难以直接统计所有人的财富、迁徙、信仰和动机。更稳妥的读法是区分诏令与实施、军报与人口、行政称谓与自我认同，并让地方材料的缺环限制解释的范围。

由此可见，该事件不是孤立的年表点，而是资源、信息与权力关系重新配置的一环。不同地区的官员、社群和家庭可能以不同方式承担征发、保护、迁移或宗教活动；现存来源只能照亮其中一部分，不能替沉默者制造统一的经验。

\eventanchor{JH-LEDGER-019-4}{永嘉元年（307年）·“司马睿出镇建业，王导等协助联络江东士族并重整东南军政”的执行中的空间与人员处置}账本记录西晋江东在永嘉元年（307年）发生““司马睿出镇建业，王导等协助联络江东士族并重整东南军政”的执行中的空间与人员处置”。条目把背景概括为“司马睿出镇建业，王导等协助联络江东士族并重整东南军政”发生的政治与区域条件共同作用，并指出该节点改变后续资源、联盟或制度处置的条件。这个节点进入区域社会时，要经由文书、道路、军队、市场、寺院和家庭等不同中介；因此，政治行动的宣布、地方接收和日常生活的改变不能被视作同一时刻完成。

本条依据《晋书·元帝纪》；《晋书·王导传》；《建康实录》，证据提示为“核心事项已有既有账本来源；本分解节点的参与者、执行范围和时间先后仍须逐案核验”。这些材料可支持年序、官方决定、战事或特定地点的线索，却难以直接统计所有人的财富、迁徙、信仰和动机。更稳妥的读法是区分诏令与实施、军报与人口、行政称谓与自我认同，并让地方材料的缺环限制解释的范围。

由此可见，该事件不是孤立的年表点，而是资源、信息与权力关系重新配置的一环。不同地区的官员、社群和家庭可能以不同方式承担征发、保护、迁移或宗教活动；现存来源只能照亮其中一部分，不能替沉默者制造统一的经验。

\eventanchor{JH-LEDGER-019-5}{永嘉元年（307年）·“司马睿出镇建业，王导等协助联络江东士族并重整东南军政”的后续制度或军政调整}账本记录西晋江东在永嘉元年（307年）发生““司马睿出镇建业，王导等协助联络江东士族并重整东南军政”的后续制度或军政调整”。条目把背景概括为核心节点发生后仍须处理其遗留问题，并指出后续影响须在不同地区和材料类型中分别核验。这个节点进入区域社会时，要经由文书、道路、军队、市场、寺院和家庭等不同中介；因此，政治行动的宣布、地方接收和日常生活的改变不能被视作同一时刻完成。

本条依据《晋书·元帝纪》；《晋书·王导传》；《建康实录》，证据提示为“核心事项已有既有账本来源；本分解节点的参与者、执行范围和时间先后仍须逐案核验”。这些材料可支持年序、官方决定、战事或特定地点的线索，却难以直接统计所有人的财富、迁徙、信仰和动机。更稳妥的读法是区分诏令与实施、军报与人口、行政称谓与自我认同，并让地方材料的缺环限制解释的范围。

由此可见，该事件不是孤立的年表点，而是资源、信息与权力关系重新配置的一环。不同地区的官员、社群和家庭可能以不同方式承担征发、保护、迁移或宗教活动；现存来源只能照亮其中一部分，不能替沉默者制造统一的经验。

\eventanchor{JH-LEDGER-019-6}{永嘉元年（307年）·“司马睿出镇建业，王导等协助联络江东士族并重整东南军政”的异源材料与影响范围的复核}账本记录西晋江东在永嘉元年（307年）发生““司马睿出镇建业，王导等协助联络江东士族并重整东南军政”的异源材料与影响范围的复核”。条目把背景概括为核心节点发生后仍须处理其遗留问题，并指出后续影响须在不同地区和材料类型中分别核验。这个节点进入区域社会时，要经由文书、道路、军队、市场、寺院和家庭等不同中介；因此，政治行动的宣布、地方接收和日常生活的改变不能被视作同一时刻完成。

本条依据《晋书·元帝纪》；《晋书·王导传》；《建康实录》，证据提示为“核心事项已有既有账本来源；本分解节点的参与者、执行范围和时间先后仍须逐案核验”。这些材料可支持年序、官方决定、战事或特定地点的线索，却难以直接统计所有人的财富、迁徙、信仰和动机。更稳妥的读法是区分诏令与实施、军报与人口、行政称谓与自我认同，并让地方材料的缺环限制解释的范围。

由此可见，该事件不是孤立的年表点，而是资源、信息与权力关系重新配置的一环。不同地区的官员、社群和家庭可能以不同方式承担征发、保护、迁移或宗教活动；现存来源只能照亮其中一部分，不能替沉默者制造统一的经验。

\eventanchor{JH-LEDGER-027-1}{建兴元年（313年）·“祖逖率部北渡长江，经屯田与结盟经营豫州、河南一带”的前期动员与资源准备}账本记录东晋前身与河南地方集团在建兴元年（313年）发生““祖逖率部北渡长江，经屯田与结盟经营豫州、河南一带”的前期动员与资源准备”。条目把背景概括为当时的权力、资源与信息约束推动相关过程展开，并指出这一过程为“祖逖率部北渡长江，经屯田与结盟经营豫州、河南一带”提供可观察的条件。这个节点进入区域社会时，要经由文书、道路、军队、市场、寺院和家庭等不同中介；因此，政治行动的宣布、地方接收和日常生活的改变不能被视作同一时刻完成。

本条依据《晋书·祖逖传》；《资治通鉴·晋纪》；《世说新语》注引《晋阳秋》，证据提示为“核心事项已有既有账本来源；本分解节点的参与者、执行范围和时间先后仍须逐案核验”。这些材料可支持年序、官方决定、战事或特定地点的线索，却难以直接统计所有人的财富、迁徙、信仰和动机。更稳妥的读法是区分诏令与实施、军报与人口、行政称谓与自我认同，并让地方材料的缺环限制解释的范围。

由此可见，该事件不是孤立的年表点，而是资源、信息与权力关系重新配置的一环。不同地区的官员、社群和家庭可能以不同方式承担征发、保护、迁移或宗教活动；现存来源只能照亮其中一部分，不能替沉默者制造统一的经验。

\eventanchor{JH-LEDGER-027-2}{建兴元年（313年）·“祖逖率部北渡长江，经屯田与结盟经营豫州、河南一带”的命令、联盟或地方响应}账本记录东晋前身与河南地方集团在建兴元年（313年）发生““祖逖率部北渡长江，经屯田与结盟经营豫州、河南一带”的命令、联盟或地方响应”。条目把背景概括为当时的权力、资源与信息约束推动相关过程展开，并指出这一过程为“祖逖率部北渡长江，经屯田与结盟经营豫州、河南一带”提供可观察的条件。这个节点进入区域社会时，要经由文书、道路、军队、市场、寺院和家庭等不同中介；因此，政治行动的宣布、地方接收和日常生活的改变不能被视作同一时刻完成。

本条依据《晋书·祖逖传》；《资治通鉴·晋纪》；《世说新语》注引《晋阳秋》，证据提示为“核心事项已有既有账本来源；本分解节点的参与者、执行范围和时间先后仍须逐案核验”。这些材料可支持年序、官方决定、战事或特定地点的线索，却难以直接统计所有人的财富、迁徙、信仰和动机。更稳妥的读法是区分诏令与实施、军报与人口、行政称谓与自我认同，并让地方材料的缺环限制解释的范围。

由此可见，该事件不是孤立的年表点，而是资源、信息与权力关系重新配置的一环。不同地区的官员、社群和家庭可能以不同方式承担征发、保护、迁移或宗教活动；现存来源只能照亮其中一部分，不能替沉默者制造统一的经验。

\eventanchor{JH-LEDGER-027-3}{建兴元年（313年）·祖逖率部北渡长江，经屯田与结盟经营豫州、河南一带}账本记录东晋前身与河南地方集团在建兴元年（313年）发生“祖逖率部北渡长江，经屯田与结盟经营豫州、河南一带”。条目把背景概括为“祖逖率部北渡长江，经屯田与结盟经营豫州、河南一带”发生的政治与区域条件共同作用，并指出该节点改变后续资源、联盟或制度处置的条件。这个节点进入区域社会时，要经由文书、道路、军队、市场、寺院和家庭等不同中介；因此，政治行动的宣布、地方接收和日常生活的改变不能被视作同一时刻完成。

本条依据《晋书·祖逖传》；《资治通鉴·晋纪》；《世说新语》注引《晋阳秋》，证据提示为“核心事项已有既有账本来源；本分解节点的参与者、执行范围和时间先后仍须逐案核验”。这些材料可支持年序、官方决定、战事或特定地点的线索，却难以直接统计所有人的财富、迁徙、信仰和动机。更稳妥的读法是区分诏令与实施、军报与人口、行政称谓与自我认同，并让地方材料的缺环限制解释的范围。

由此可见，该事件不是孤立的年表点，而是资源、信息与权力关系重新配置的一环。不同地区的官员、社群和家庭可能以不同方式承担征发、保护、迁移或宗教活动；现存来源只能照亮其中一部分，不能替沉默者制造统一的经验。

\eventanchor{JH-LEDGER-027-4}{建兴元年（313年）·“祖逖率部北渡长江，经屯田与结盟经营豫州、河南一带”的执行中的空间与人员处置}账本记录东晋前身与河南地方集团在建兴元年（313年）发生““祖逖率部北渡长江，经屯田与结盟经营豫州、河南一带”的执行中的空间与人员处置”。条目把背景概括为“祖逖率部北渡长江，经屯田与结盟经营豫州、河南一带”发生的政治与区域条件共同作用，并指出该节点改变后续资源、联盟或制度处置的条件。这个节点进入区域社会时，要经由文书、道路、军队、市场、寺院和家庭等不同中介；因此，政治行动的宣布、地方接收和日常生活的改变不能被视作同一时刻完成。

本条依据《晋书·祖逖传》；《资治通鉴·晋纪》；《世说新语》注引《晋阳秋》，证据提示为“核心事项已有既有账本来源；本分解节点的参与者、执行范围和时间先后仍须逐案核验”。这些材料可支持年序、官方决定、战事或特定地点的线索，却难以直接统计所有人的财富、迁徙、信仰和动机。更稳妥的读法是区分诏令与实施、军报与人口、行政称谓与自我认同，并让地方材料的缺环限制解释的范围。

由此可见，该事件不是孤立的年表点，而是资源、信息与权力关系重新配置的一环。不同地区的官员、社群和家庭可能以不同方式承担征发、保护、迁移或宗教活动；现存来源只能照亮其中一部分，不能替沉默者制造统一的经验。

\eventanchor{JH-LEDGER-027-5}{建兴元年（313年）·“祖逖率部北渡长江，经屯田与结盟经营豫州、河南一带”的后续制度或军政调整}账本记录东晋前身与河南地方集团在建兴元年（313年）发生““祖逖率部北渡长江，经屯田与结盟经营豫州、河南一带”的后续制度或军政调整”。条目把背景概括为核心节点发生后仍须处理其遗留问题，并指出后续影响须在不同地区和材料类型中分别核验。这个节点进入区域社会时，要经由文书、道路、军队、市场、寺院和家庭等不同中介；因此，政治行动的宣布、地方接收和日常生活的改变不能被视作同一时刻完成。

本条依据《晋书·祖逖传》；《资治通鉴·晋纪》；《世说新语》注引《晋阳秋》，证据提示为“核心事项已有既有账本来源；本分解节点的参与者、执行范围和时间先后仍须逐案核验”。这些材料可支持年序、官方决定、战事或特定地点的线索，却难以直接统计所有人的财富、迁徙、信仰和动机。更稳妥的读法是区分诏令与实施、军报与人口、行政称谓与自我认同，并让地方材料的缺环限制解释的范围。

由此可见，该事件不是孤立的年表点，而是资源、信息与权力关系重新配置的一环。不同地区的官员、社群和家庭可能以不同方式承担征发、保护、迁移或宗教活动；现存来源只能照亮其中一部分，不能替沉默者制造统一的经验。

\eventanchor{JH-LEDGER-027-6}{建兴元年（313年）·“祖逖率部北渡长江，经屯田与结盟经营豫州、河南一带”的异源材料与影响范围的复核}账本记录东晋前身与河南地方集团在建兴元年（313年）发生““祖逖率部北渡长江，经屯田与结盟经营豫州、河南一带”的异源材料与影响范围的复核”。条目把背景概括为核心节点发生后仍须处理其遗留问题，并指出后续影响须在不同地区和材料类型中分别核验。这个节点进入区域社会时，要经由文书、道路、军队、市场、寺院和家庭等不同中介；因此，政治行动的宣布、地方接收和日常生活的改变不能被视作同一时刻完成。

本条依据《晋书·祖逖传》；《资治通鉴·晋纪》；《世说新语》注引《晋阳秋》，证据提示为“核心事项已有既有账本来源；本分解节点的参与者、执行范围和时间先后仍须逐案核验”。这些材料可支持年序、官方决定、战事或特定地点的线索，却难以直接统计所有人的财富、迁徙、信仰和动机。更稳妥的读法是区分诏令与实施、军报与人口、行政称谓与自我认同，并让地方材料的缺环限制解释的范围。

由此可见，该事件不是孤立的年表点，而是资源、信息与权力关系重新配置的一环。不同地区的官员、社群和家庭可能以不同方式承担征发、保护、迁移或宗教活动；现存来源只能照亮其中一部分，不能替沉默者制造统一的经验。

\eventanchor{JH-LEDGER-028-1}{建兴三年（315年）·“拓跋猗卢受晋封为代王，在北部边地扩展部众与城郭控制”的前期动员与资源准备}账本记录代国与西晋残余政权在建兴三年（315年）发生““拓跋猗卢受晋封为代王，在北部边地扩展部众与城郭控制”的前期动员与资源准备”。条目把背景概括为当时的权力、资源与信息约束推动相关过程展开，并指出这一过程为“拓跋猗卢受晋封为代王，在北部边地扩展部众与城郭控制”提供可观察的条件。这个节点进入区域社会时，要经由文书、道路、军队、市场、寺院和家庭等不同中介；因此，政治行动的宣布、地方接收和日常生活的改变不能被视作同一时刻完成。

本条依据《魏书·序纪》；《晋书·愍帝纪》；《资治通鉴·晋纪》，证据提示为“核心事项已有既有账本来源；本分解节点的参与者、执行范围和时间先后仍须逐案核验”。这些材料可支持年序、官方决定、战事或特定地点的线索，却难以直接统计所有人的财富、迁徙、信仰和动机。更稳妥的读法是区分诏令与实施、军报与人口、行政称谓与自我认同，并让地方材料的缺环限制解释的范围。

由此可见，该事件不是孤立的年表点，而是资源、信息与权力关系重新配置的一环。不同地区的官员、社群和家庭可能以不同方式承担征发、保护、迁移或宗教活动；现存来源只能照亮其中一部分，不能替沉默者制造统一的经验。

\eventanchor{JH-LEDGER-028-2}{建兴三年（315年）·“拓跋猗卢受晋封为代王，在北部边地扩展部众与城郭控制”的命令、联盟或地方响应}账本记录代国与西晋残余政权在建兴三年（315年）发生““拓跋猗卢受晋封为代王，在北部边地扩展部众与城郭控制”的命令、联盟或地方响应”。条目把背景概括为当时的权力、资源与信息约束推动相关过程展开，并指出这一过程为“拓跋猗卢受晋封为代王，在北部边地扩展部众与城郭控制”提供可观察的条件。这个节点进入区域社会时，要经由文书、道路、军队、市场、寺院和家庭等不同中介；因此，政治行动的宣布、地方接收和日常生活的改变不能被视作同一时刻完成。

本条依据《魏书·序纪》；《晋书·愍帝纪》；《资治通鉴·晋纪》，证据提示为“核心事项已有既有账本来源；本分解节点的参与者、执行范围和时间先后仍须逐案核验”。这些材料可支持年序、官方决定、战事或特定地点的线索，却难以直接统计所有人的财富、迁徙、信仰和动机。更稳妥的读法是区分诏令与实施、军报与人口、行政称谓与自我认同，并让地方材料的缺环限制解释的范围。

由此可见，该事件不是孤立的年表点，而是资源、信息与权力关系重新配置的一环。不同地区的官员、社群和家庭可能以不同方式承担征发、保护、迁移或宗教活动；现存来源只能照亮其中一部分，不能替沉默者制造统一的经验。

\eventanchor{JH-LEDGER-028-3}{建兴三年（315年）·拓跋猗卢受晋封为代王，在北部边地扩展部众与城郭控制}账本记录代国与西晋残余政权在建兴三年（315年）发生“拓跋猗卢受晋封为代王，在北部边地扩展部众与城郭控制”。条目把背景概括为“拓跋猗卢受晋封为代王，在北部边地扩展部众与城郭控制”发生的政治与区域条件共同作用，并指出该节点改变后续资源、联盟或制度处置的条件。这个节点进入区域社会时，要经由文书、道路、军队、市场、寺院和家庭等不同中介；因此，政治行动的宣布、地方接收和日常生活的改变不能被视作同一时刻完成。

本条依据《魏书·序纪》；《晋书·愍帝纪》；《资治通鉴·晋纪》，证据提示为“核心事项已有既有账本来源；本分解节点的参与者、执行范围和时间先后仍须逐案核验”。这些材料可支持年序、官方决定、战事或特定地点的线索，却难以直接统计所有人的财富、迁徙、信仰和动机。更稳妥的读法是区分诏令与实施、军报与人口、行政称谓与自我认同，并让地方材料的缺环限制解释的范围。

由此可见，该事件不是孤立的年表点，而是资源、信息与权力关系重新配置的一环。不同地区的官员、社群和家庭可能以不同方式承担征发、保护、迁移或宗教活动；现存来源只能照亮其中一部分，不能替沉默者制造统一的经验。

\eventanchor{JH-LEDGER-028-4}{建兴三年（315年）·“拓跋猗卢受晋封为代王，在北部边地扩展部众与城郭控制”的执行中的空间与人员处置}账本记录代国与西晋残余政权在建兴三年（315年）发生““拓跋猗卢受晋封为代王，在北部边地扩展部众与城郭控制”的执行中的空间与人员处置”。条目把背景概括为“拓跋猗卢受晋封为代王，在北部边地扩展部众与城郭控制”发生的政治与区域条件共同作用，并指出该节点改变后续资源、联盟或制度处置的条件。这个节点进入区域社会时，要经由文书、道路、军队、市场、寺院和家庭等不同中介；因此，政治行动的宣布、地方接收和日常生活的改变不能被视作同一时刻完成。

本条依据《魏书·序纪》；《晋书·愍帝纪》；《资治通鉴·晋纪》，证据提示为“核心事项已有既有账本来源；本分解节点的参与者、执行范围和时间先后仍须逐案核验”。这些材料可支持年序、官方决定、战事或特定地点的线索，却难以直接统计所有人的财富、迁徙、信仰和动机。更稳妥的读法是区分诏令与实施、军报与人口、行政称谓与自我认同，并让地方材料的缺环限制解释的范围。

由此可见，该事件不是孤立的年表点，而是资源、信息与权力关系重新配置的一环。不同地区的官员、社群和家庭可能以不同方式承担征发、保护、迁移或宗教活动；现存来源只能照亮其中一部分，不能替沉默者制造统一的经验。

\eventanchor{JH-LEDGER-028-5}{建兴三年（315年）·“拓跋猗卢受晋封为代王，在北部边地扩展部众与城郭控制”的后续制度或军政调整}账本记录代国与西晋残余政权在建兴三年（315年）发生““拓跋猗卢受晋封为代王，在北部边地扩展部众与城郭控制”的后续制度或军政调整”。条目把背景概括为核心节点发生后仍须处理其遗留问题，并指出后续影响须在不同地区和材料类型中分别核验。这个节点进入区域社会时，要经由文书、道路、军队、市场、寺院和家庭等不同中介；因此，政治行动的宣布、地方接收和日常生活的改变不能被视作同一时刻完成。

本条依据《魏书·序纪》；《晋书·愍帝纪》；《资治通鉴·晋纪》，证据提示为“核心事项已有既有账本来源；本分解节点的参与者、执行范围和时间先后仍须逐案核验”。这些材料可支持年序、官方决定、战事或特定地点的线索，却难以直接统计所有人的财富、迁徙、信仰和动机。更稳妥的读法是区分诏令与实施、军报与人口、行政称谓与自我认同，并让地方材料的缺环限制解释的范围。

由此可见，该事件不是孤立的年表点，而是资源、信息与权力关系重新配置的一环。不同地区的官员、社群和家庭可能以不同方式承担征发、保护、迁移或宗教活动；现存来源只能照亮其中一部分，不能替沉默者制造统一的经验。

\eventanchor{JH-LEDGER-028-6}{建兴三年（315年）·“拓跋猗卢受晋封为代王，在北部边地扩展部众与城郭控制”的异源材料与影响范围的复核}账本记录代国与西晋残余政权在建兴三年（315年）发生““拓跋猗卢受晋封为代王，在北部边地扩展部众与城郭控制”的异源材料与影响范围的复核”。条目把背景概括为核心节点发生后仍须处理其遗留问题，并指出后续影响须在不同地区和材料类型中分别核验。这个节点进入区域社会时，要经由文书、道路、军队、市场、寺院和家庭等不同中介；因此，政治行动的宣布、地方接收和日常生活的改变不能被视作同一时刻完成。

本条依据《魏书·序纪》；《晋书·愍帝纪》；《资治通鉴·晋纪》，证据提示为“核心事项已有既有账本来源；本分解节点的参与者、执行范围和时间先后仍须逐案核验”。这些材料可支持年序、官方决定、战事或特定地点的线索，却难以直接统计所有人的财富、迁徙、信仰和动机。更稳妥的读法是区分诏令与实施、军报与人口、行政称谓与自我认同，并让地方材料的缺环限制解释的范围。

由此可见，该事件不是孤立的年表点，而是资源、信息与权力关系重新配置的一环。不同地区的官员、社群和家庭可能以不同方式承担征发、保护、迁移或宗教活动；现存来源只能照亮其中一部分，不能替沉默者制造统一的经验。

\eventanchor{JH-LEDGER-030-1}{建武元年三月（317年）·“司马睿在建康即晋王位，建立承继西晋的南方朝廷框架”的前期动员与资源准备}账本记录东晋在建武元年三月（317年）发生““司马睿在建康即晋王位，建立承继西晋的南方朝廷框架”的前期动员与资源准备”。条目把背景概括为当时的权力、资源与信息约束推动相关过程展开，并指出这一过程为“司马睿在建康即晋王位，建立承继西晋的南方朝廷框架”提供可观察的条件。这个节点进入区域社会时，要经由文书、道路、军队、市场、寺院和家庭等不同中介；因此，政治行动的宣布、地方接收和日常生活的改变不能被视作同一时刻完成。

本条依据《晋书·元帝纪》；《晋书·王导传》；《建康实录》，证据提示为“核心事项已有既有账本来源；本分解节点的参与者、执行范围和时间先后仍须逐案核验”。这些材料可支持年序、官方决定、战事或特定地点的线索，却难以直接统计所有人的财富、迁徙、信仰和动机。更稳妥的读法是区分诏令与实施、军报与人口、行政称谓与自我认同，并让地方材料的缺环限制解释的范围。

由此可见，该事件不是孤立的年表点，而是资源、信息与权力关系重新配置的一环。不同地区的官员、社群和家庭可能以不同方式承担征发、保护、迁移或宗教活动；现存来源只能照亮其中一部分，不能替沉默者制造统一的经验。

\eventanchor{JH-LEDGER-030-2}{建武元年三月（317年）·“司马睿在建康即晋王位，建立承继西晋的南方朝廷框架”的命令、联盟或地方响应}账本记录东晋在建武元年三月（317年）发生““司马睿在建康即晋王位，建立承继西晋的南方朝廷框架”的命令、联盟或地方响应”。条目把背景概括为当时的权力、资源与信息约束推动相关过程展开，并指出这一过程为“司马睿在建康即晋王位，建立承继西晋的南方朝廷框架”提供可观察的条件。这个节点进入区域社会时，要经由文书、道路、军队、市场、寺院和家庭等不同中介；因此，政治行动的宣布、地方接收和日常生活的改变不能被视作同一时刻完成。

本条依据《晋书·元帝纪》；《晋书·王导传》；《建康实录》，证据提示为“核心事项已有既有账本来源；本分解节点的参与者、执行范围和时间先后仍须逐案核验”。这些材料可支持年序、官方决定、战事或特定地点的线索，却难以直接统计所有人的财富、迁徙、信仰和动机。更稳妥的读法是区分诏令与实施、军报与人口、行政称谓与自我认同，并让地方材料的缺环限制解释的范围。

由此可见，该事件不是孤立的年表点，而是资源、信息与权力关系重新配置的一环。不同地区的官员、社群和家庭可能以不同方式承担征发、保护、迁移或宗教活动；现存来源只能照亮其中一部分，不能替沉默者制造统一的经验。

\eventanchor{JH-LEDGER-030-3}{建武元年三月（317年）·司马睿在建康即晋王位，建立承继西晋的南方朝廷框架}账本记录东晋在建武元年三月（317年）发生“司马睿在建康即晋王位，建立承继西晋的南方朝廷框架”。条目把背景概括为“司马睿在建康即晋王位，建立承继西晋的南方朝廷框架”发生的政治与区域条件共同作用，并指出该节点改变后续资源、联盟或制度处置的条件。这个节点进入区域社会时，要经由文书、道路、军队、市场、寺院和家庭等不同中介；因此，政治行动的宣布、地方接收和日常生活的改变不能被视作同一时刻完成。

本条依据《晋书·元帝纪》；《晋书·王导传》；《建康实录》，证据提示为“核心事项已有既有账本来源；本分解节点的参与者、执行范围和时间先后仍须逐案核验”。这些材料可支持年序、官方决定、战事或特定地点的线索，却难以直接统计所有人的财富、迁徙、信仰和动机。更稳妥的读法是区分诏令与实施、军报与人口、行政称谓与自我认同，并让地方材料的缺环限制解释的范围。

由此可见，该事件不是孤立的年表点，而是资源、信息与权力关系重新配置的一环。不同地区的官员、社群和家庭可能以不同方式承担征发、保护、迁移或宗教活动；现存来源只能照亮其中一部分，不能替沉默者制造统一的经验。

\eventanchor{JH-LEDGER-030-4}{建武元年三月（317年）·“司马睿在建康即晋王位，建立承继西晋的南方朝廷框架”的执行中的空间与人员处置}账本记录东晋在建武元年三月（317年）发生““司马睿在建康即晋王位，建立承继西晋的南方朝廷框架”的执行中的空间与人员处置”。条目把背景概括为“司马睿在建康即晋王位，建立承继西晋的南方朝廷框架”发生的政治与区域条件共同作用，并指出该节点改变后续资源、联盟或制度处置的条件。这个节点进入区域社会时，要经由文书、道路、军队、市场、寺院和家庭等不同中介；因此，政治行动的宣布、地方接收和日常生活的改变不能被视作同一时刻完成。

本条依据《晋书·元帝纪》；《晋书·王导传》；《建康实录》，证据提示为“核心事项已有既有账本来源；本分解节点的参与者、执行范围和时间先后仍须逐案核验”。这些材料可支持年序、官方决定、战事或特定地点的线索，却难以直接统计所有人的财富、迁徙、信仰和动机。更稳妥的读法是区分诏令与实施、军报与人口、行政称谓与自我认同，并让地方材料的缺环限制解释的范围。

由此可见，该事件不是孤立的年表点，而是资源、信息与权力关系重新配置的一环。不同地区的官员、社群和家庭可能以不同方式承担征发、保护、迁移或宗教活动；现存来源只能照亮其中一部分，不能替沉默者制造统一的经验。

\eventanchor{JH-LEDGER-030-5}{建武元年三月（317年）·“司马睿在建康即晋王位，建立承继西晋的南方朝廷框架”的后续制度或军政调整}账本记录东晋在建武元年三月（317年）发生““司马睿在建康即晋王位，建立承继西晋的南方朝廷框架”的后续制度或军政调整”。条目把背景概括为核心节点发生后仍须处理其遗留问题，并指出后续影响须在不同地区和材料类型中分别核验。这个节点进入区域社会时，要经由文书、道路、军队、市场、寺院和家庭等不同中介；因此，政治行动的宣布、地方接收和日常生活的改变不能被视作同一时刻完成。

本条依据《晋书·元帝纪》；《晋书·王导传》；《建康实录》，证据提示为“核心事项已有既有账本来源；本分解节点的参与者、执行范围和时间先后仍须逐案核验”。这些材料可支持年序、官方决定、战事或特定地点的线索，却难以直接统计所有人的财富、迁徙、信仰和动机。更稳妥的读法是区分诏令与实施、军报与人口、行政称谓与自我认同，并让地方材料的缺环限制解释的范围。

由此可见，该事件不是孤立的年表点，而是资源、信息与权力关系重新配置的一环。不同地区的官员、社群和家庭可能以不同方式承担征发、保护、迁移或宗教活动；现存来源只能照亮其中一部分，不能替沉默者制造统一的经验。

\eventanchor{JH-LEDGER-030-6}{建武元年三月（317年）·“司马睿在建康即晋王位，建立承继西晋的南方朝廷框架”的异源材料与影响范围的复核}账本记录东晋在建武元年三月（317年）发生““司马睿在建康即晋王位，建立承继西晋的南方朝廷框架”的异源材料与影响范围的复核”。条目把背景概括为核心节点发生后仍须处理其遗留问题，并指出后续影响须在不同地区和材料类型中分别核验。这个节点进入区域社会时，要经由文书、道路、军队、市场、寺院和家庭等不同中介；因此，政治行动的宣布、地方接收和日常生活的改变不能被视作同一时刻完成。

本条依据《晋书·元帝纪》；《晋书·王导传》；《建康实录》，证据提示为“核心事项已有既有账本来源；本分解节点的参与者、执行范围和时间先后仍须逐案核验”。这些材料可支持年序、官方决定、战事或特定地点的线索，却难以直接统计所有人的财富、迁徙、信仰和动机。更稳妥的读法是区分诏令与实施、军报与人口、行政称谓与自我认同，并让地方材料的缺环限制解释的范围。

由此可见，该事件不是孤立的年表点，而是资源、信息与权力关系重新配置的一环。不同地区的官员、社群和家庭可能以不同方式承担征发、保护、迁移或宗教活动；现存来源只能照亮其中一部分，不能替沉默者制造统一的经验。

\eventanchor{JH-LEDGER-031-1}{太兴元年三月（318年）·“司马睿称帝，是为元帝，东晋以建康为都正式确立”的前期动员与资源准备}账本记录东晋在太兴元年三月（318年）发生““司马睿称帝，是为元帝，东晋以建康为都正式确立”的前期动员与资源准备”。条目把背景概括为当时的权力、资源与信息约束推动相关过程展开，并指出这一过程为“司马睿称帝，是为元帝，东晋以建康为都正式确立”提供可观察的条件。这个节点进入区域社会时，要经由文书、道路、军队、市场、寺院和家庭等不同中介；因此，政治行动的宣布、地方接收和日常生活的改变不能被视作同一时刻完成。

本条依据《晋书·元帝纪》；《晋书·王导传》；《资治通鉴·晋纪》；《建康实录》，证据提示为“核心事项已有既有账本来源；本分解节点的参与者、执行范围和时间先后仍须逐案核验”。这些材料可支持年序、官方决定、战事或特定地点的线索，却难以直接统计所有人的财富、迁徙、信仰和动机。更稳妥的读法是区分诏令与实施、军报与人口、行政称谓与自我认同，并让地方材料的缺环限制解释的范围。

由此可见，该事件不是孤立的年表点，而是资源、信息与权力关系重新配置的一环。不同地区的官员、社群和家庭可能以不同方式承担征发、保护、迁移或宗教活动；现存来源只能照亮其中一部分，不能替沉默者制造统一的经验。

\eventanchor{JH-LEDGER-031-2}{太兴元年三月（318年）·“司马睿称帝，是为元帝，东晋以建康为都正式确立”的命令、联盟或地方响应}账本记录东晋在太兴元年三月（318年）发生““司马睿称帝，是为元帝，东晋以建康为都正式确立”的命令、联盟或地方响应”。条目把背景概括为当时的权力、资源与信息约束推动相关过程展开，并指出这一过程为“司马睿称帝，是为元帝，东晋以建康为都正式确立”提供可观察的条件。这个节点进入区域社会时，要经由文书、道路、军队、市场、寺院和家庭等不同中介；因此，政治行动的宣布、地方接收和日常生活的改变不能被视作同一时刻完成。

本条依据《晋书·元帝纪》；《晋书·王导传》；《资治通鉴·晋纪》；《建康实录》，证据提示为“核心事项已有既有账本来源；本分解节点的参与者、执行范围和时间先后仍须逐案核验”。这些材料可支持年序、官方决定、战事或特定地点的线索，却难以直接统计所有人的财富、迁徙、信仰和动机。更稳妥的读法是区分诏令与实施、军报与人口、行政称谓与自我认同，并让地方材料的缺环限制解释的范围。

由此可见，该事件不是孤立的年表点，而是资源、信息与权力关系重新配置的一环。不同地区的官员、社群和家庭可能以不同方式承担征发、保护、迁移或宗教活动；现存来源只能照亮其中一部分，不能替沉默者制造统一的经验。

\eventanchor{JH-LEDGER-031-3}{太兴元年三月（318年）·司马睿称帝，是为元帝，东晋以建康为都正式确立}账本记录东晋在太兴元年三月（318年）发生“司马睿称帝，是为元帝，东晋以建康为都正式确立”。条目把背景概括为“司马睿称帝，是为元帝，东晋以建康为都正式确立”发生的政治与区域条件共同作用，并指出该节点改变后续资源、联盟或制度处置的条件。这个节点进入区域社会时，要经由文书、道路、军队、市场、寺院和家庭等不同中介；因此，政治行动的宣布、地方接收和日常生活的改变不能被视作同一时刻完成。

本条依据《晋书·元帝纪》；《晋书·王导传》；《资治通鉴·晋纪》；《建康实录》，证据提示为“核心事项已有既有账本来源；本分解节点的参与者、执行范围和时间先后仍须逐案核验”。这些材料可支持年序、官方决定、战事或特定地点的线索，却难以直接统计所有人的财富、迁徙、信仰和动机。更稳妥的读法是区分诏令与实施、军报与人口、行政称谓与自我认同，并让地方材料的缺环限制解释的范围。

由此可见，该事件不是孤立的年表点，而是资源、信息与权力关系重新配置的一环。不同地区的官员、社群和家庭可能以不同方式承担征发、保护、迁移或宗教活动；现存来源只能照亮其中一部分，不能替沉默者制造统一的经验。

\eventanchor{JH-LEDGER-031-4}{太兴元年三月（318年）·“司马睿称帝，是为元帝，东晋以建康为都正式确立”的执行中的空间与人员处置}账本记录东晋在太兴元年三月（318年）发生““司马睿称帝，是为元帝，东晋以建康为都正式确立”的执行中的空间与人员处置”。条目把背景概括为“司马睿称帝，是为元帝，东晋以建康为都正式确立”发生的政治与区域条件共同作用，并指出该节点改变后续资源、联盟或制度处置的条件。这个节点进入区域社会时，要经由文书、道路、军队、市场、寺院和家庭等不同中介；因此，政治行动的宣布、地方接收和日常生活的改变不能被视作同一时刻完成。

本条依据《晋书·元帝纪》；《晋书·王导传》；《资治通鉴·晋纪》；《建康实录》，证据提示为“核心事项已有既有账本来源；本分解节点的参与者、执行范围和时间先后仍须逐案核验”。这些材料可支持年序、官方决定、战事或特定地点的线索，却难以直接统计所有人的财富、迁徙、信仰和动机。更稳妥的读法是区分诏令与实施、军报与人口、行政称谓与自我认同，并让地方材料的缺环限制解释的范围。

由此可见，该事件不是孤立的年表点，而是资源、信息与权力关系重新配置的一环。不同地区的官员、社群和家庭可能以不同方式承担征发、保护、迁移或宗教活动；现存来源只能照亮其中一部分，不能替沉默者制造统一的经验。
